\documentclass[12pt,a4paper,UTF8]{ctexart}
\usepackage[a4paper,top=22mm,bottom=22mm,left=22mm,right=22mm,headheight=15pt]{geometry}
\usepackage{amsmath,amssymb,mathtools,bm}
\usepackage{graphicx,xcolor,booktabs,longtable,array,enumitem,fancyhdr,caption,float,microtype,setspace}
\newcommand{\Needspace}[1]{\par\begingroup\dimen0=#1\relax\dimen2=\pagegoal\advance\dimen2 by-\pagetotal\ifdim\dimen0>\dimen2\newpage\fi\endgroup}
\newcommand{\needspace}[1]{\par\begingroup\skip0=\lastskip\vskip-\skip0\dimen0=#1\relax\advance\dimen0 by\skip0\vskip\dimen0\penalty-100\vskip-\dimen0\vskip\skip0\penalty9999\endgroup}
\usepackage[hidelinks,unicode]{hyperref}
\usepackage[most]{tcolorbox}
\definecolor{ContestBlue}{HTML}{17365D}
\hypersetup{pdftitle={第 043 套},pdfauthor={}}
\setstretch{1.24}
\setlength{\parindent}{2em}
\setlength{\parskip}{0.35em}
\setlength{\emergencystretch}{2em}
\setlist{itemsep=0.35em,topsep=0.4em}
\captionsetup{font=small,labelfont=bf,justification=centering}
\setcounter{secnumdepth}{0}
\pagestyle{fancy}
\fancyhf{}
\fancyhead[L]{\small 第 043 套}
\fancyfoot[C]{\thepage}
\renewcommand{\headrulewidth}{0.4pt}
\ctexset{
  section={format=\Large\bfseries\color{ContestBlue},beforeskip=1.2em,afterskip=0.7em},
  subsection={format=\large\bfseries\color{ContestBlue},beforeskip=1.0em,afterskip=0.5em},
  subsubsection={format=\normalsize\bfseries,beforeskip=0.8em,afterskip=0.35em}
}
\makeatletter
\newcommand{\tightlist}{\setlength{\itemsep}{0pt}\setlength{\parskip}{0pt}}
\newcommand{\maxwidth}{\ifdim\Gin@nat@width>\linewidth\linewidth\else\Gin@nat@width\fi}
\newcommand{\maxheight}{\ifdim\Gin@nat@height>0.72\textheight 0.72\textheight\else\Gin@nat@height\fi}
\newcommand{\pandocbounded}[1]{\begingroup\setkeys{Gin}{width=\maxwidth,height=\maxheight,keepaspectratio}#1\endgroup}
\makeatother
\begin{document}
\begin{titlepage}
  \centering
  \vspace*{58mm}
  {\Huge\bfseries\color{ContestBlue} 第 043 套\par}
  \vfill
  \vspace*{28mm}
\end{titlepage}
\section{第 1 题}\label{ux7b2c-1-ux9898}

一列汽车排队等红灯，初始车间距为
\(2\,\mathrm m\)，这里的间距指前车尾到后车头的距离。汽车启动后以
\(1\,\mathrm{m/s^2}\) 加速；后车等前车把间距拉到 \(7\,\mathrm m\)
时才启动。每辆车加速到 \(10\,\mathrm{m/s}\) 后改为匀速。

（1）车队全部匀速行驶时，车间距是多少？

（2）前车以 \(1\,\mathrm{m/s^2}\)
减速，后车延迟一段时间后以相同加速度减速。为使此后两车间距始终大于
\(2\,\mathrm m\)，后车应在多少时间内开始减速？

\includegraphics[width=0.86\linewidth,height=\textheight,keepaspectratio,alt={题1　启动延迟与前尾---后头间距}]{../figures/question-01.pdf}

\subsection{一、题型定位}\label{ux4e00ux9898ux578bux5b9aux4f4d}

前后两车采用相同的速度程序，只是开始时刻不同。启动时用位移确定延迟，停车时用速度差的积分确定间距损失。两段运动都可以归结为同一条速度曲线的时间平移。

\subsection{二、知识点}\label{ux4e8cux77e5ux8bc6ux70b9}

匀加速位移 \(s=at^2/2\)；分段速度函数；相对速度
\(\dot d=v_F-v_R\)；速度---时间图下的面积；严格不等式与临界状态。

\subsection{三、物理图像}\label{ux4e09ux7269ux7406ux56feux50cf}

图1-1把前车尾与后车头分别作为位置标记。车长始终包含在两个标记的初始位置中，后续计算只涉及它们各自的位移。图1-2的两条速度曲线形状相同，后车曲线向右平移；阴影面积就是车距的变化量。

\subsection{四、解题入口}\label{ux56dbux89e3ux9898ux5165ux53e3}

先求后车启动前，前车多走的
\(7-2=5\,\mathrm m\)。得到启动延迟后，不必分别累加每辆车的全部位移；比较两车相同速度程序在时间上错开的部分即可。

\subsection{五、解析}\label{ux4e94ux89e3ux6790}

\textbf{启动延迟。} 记初始间距 \(d_0=2\,\mathrm m\)，启动门槛
\(d_s=7\,\mathrm m\)，加速度 \(a=1\,\mathrm{m/s^2}\)，限速
\(v=10\,\mathrm{m/s}\)。以前车启动为 \(t=0\)，后车在 \(t=\Delta\)
时启动。在这一段后车不动，故 \[
d_s=d_0+\frac12a\Delta^2.
\] 所以 \[
\boxed{\Delta=\sqrt{\frac{2(d_s-d_0)}a}=\sqrt{10}\,\mathrm s
\simeq3.16228\,\mathrm s}.
\] 前车加速到限速需要 \[
t_a=\frac va=10\,\mathrm s.
\] 由于
\(\Delta<t_a\)，上述门槛确实发生在前车的加速阶段，所用位移式成立。

\textbf{车队匀速时的间距。} 用从各车启动时刻算起的位移函数 \[
s(t)=
\begin{cases}
0,&t\le0,\\
\dfrac12at^2,&0<t<t_a,\\
vt-\dfrac{v^2}{2a},&t\ge t_a.
\end{cases}
\] 两车的间距为 \[
d(t)=d_0+s(t)-s(t-\Delta).
\] 待后车也达到限速，即 \(t\ge t_a+\Delta\)，两项都在第三段，常量
\(v^2/(2a)\) 相消： \[
\begin{aligned}
d_\infty
&=d_0+\left(vt-\frac{v^2}{2a}\right)
-\left[v(t-\Delta)-\frac{v^2}{2a}\right]\\
&=d_0+v\Delta.
\end{aligned}
\] 因此 \[
\boxed{d_\infty=2+10\sqrt{10}\simeq33.6228\,\mathrm m}.
\] \(7\,\mathrm m\)
只是后车启动时的距离。后车启动后仍比前车慢，间距继续扩大，直到两车都进入匀速阶段才固定下来。

\begin{figure}[H]
\centering
\includegraphics[width=0.88\linewidth,height=\textheight,keepaspectratio,alt={图1-2　速度曲线平移与启动、制动的面积账}]{../figures/solution-01.pdf}
\caption*{图1-2　速度曲线平移与启动、制动的面积账}
\end{figure}

\textbf{启动全过程的分段检查。} 在 \(\Delta<t<t_a\)，两车都加速， \[
v_F=at,\qquad v_R=a(t-\Delta),
\] 所以 \[
\dot d=a\Delta>0,
\] \[
d(t)=d_0+a\Delta t-\frac12a\Delta^2.
\] 在 \(t_a<t<t_a+\Delta\)，前车已匀速而后车仍加速： \[
\dot d=v-a(t-\Delta)=a(t_a+\Delta-t)>0.
\] 此段的相对速度从 \(a\Delta\)
线性降到零。整个启动过程中，间距先按二次函数增大，再匀速增大，最后以逐渐减小的增长率趋于
\(d_\infty\)；没有中途缩小的阶段。

\textbf{制动造成的距离损失。} 现把前车开始制动重新记为
\(t=0\)。两车当时都以速度 \(v\) 行驶，间距为 \(d_\infty\)。设后车延迟
\(\tau\ge0\) 后开始制动，制动加速度大小为
\(b=1\,\mathrm{m/s^2}\)。前车从开始制动到停车所走距离 \[
s_F=\frac{v^2}{2b}.
\] 后车先匀速走 \(v\tau\)，随后采用相同制动程序，因此 \[
s_R=v\tau+\frac{v^2}{2b}.
\] 两车都停车时，间距为 \[
\boxed{d_{\rm stop}=d_\infty-v\tau}.
\] 同样的制动距离相消，剩下的损失恰好是后车等待期间的 \(v\tau\)。

\textbf{为什么最终间距就是最小间距。}
从前车开始制动到后车停车，前车的速度始终不大于后车的速度。对
\(0\le\tau\le v/b\)，分三段： \[
v_F-v_R=
\begin{cases}
-bt,&0<t<\tau,\\
-b\tau,&\tau<t<v/b,\\
-v+b(t-\tau),&v/b<t<v/b+\tau.
\end{cases}
\] 每段均非正，因此 \(d(t)\) 单调减小，最后一段末尾才停止变化。若
\(\tau>v/b\)，前车先停车，后车尚未制动，间距仍只会进一步缩小，结论不变。故只需检查停车后的距离，无须额外猜测某个更危险的中间时刻。

原题要求严格大于 \(d_0\)： \[
d_\infty-v\tau>d_0.
\] 代入 \(d_\infty=d_0+v\Delta\)，得到 \[
\boxed{0\le\tau<\Delta=\sqrt{10}\,\mathrm s}.
\] 也就是后车必须在前车开始制动后的 \(3.16\,\mathrm s\)
之前启动制动。\(\tau=\sqrt{10}\,\mathrm s\) 时，最后恰好剩
\(2\,\mathrm m\)，不满足``大于''。

\textbf{按速度图独立复算。}
图1-2中启动的两条曲线相差一个时间平移。把它们之间的面积分成加速起始三角形、共同加速阶段的矩形和末尾三角形：
\[
\Delta d
=\frac12a\Delta^2
+a\Delta(t_a-\Delta)
+\frac12a\Delta^2
=a\Delta t_a=v\Delta.
\] 制动的三块面积分别是 \[
\frac12b\tau^2,\qquad
b\tau\left(\frac vb-\tau\right),\qquad
\frac12b\tau^2,
\] 其和为
\(v\tau\)。同一面积结构解释了启动拉开多少与延迟制动追回多少，而不依赖长位移表。

\textbf{不必追踪整个车队。}
每一对相邻车辆都采用相同的启动规则。当前车从静止开始运动后，后车等待它多走
\(5\,\mathrm m\)，因此每一对相邻车辆都有相同的 \(\Delta\)。第 \(j\)
辆车的启动时刻为 \((j-1)\Delta\)，但相邻车辆达到限速后的间距始终为
\(d_0+v\Delta\)，不会随车队位置累加成更大的间距。

\textbf{一般制动速度程序。} 若两车采用同一条单调下降的速度曲线
\(v_b(t)\)，前车在 \(0\) 开始，后车在 \(\tau\) 开始，且等待时保持初速度
\(v\)，则后车比前车多走的距离仍为
\(v\tau\)。这个结果来自时间平移后的积分，既不要求制动加速度恒定，也不要求其等于启动加速度。它要求的是两车从开始制动到停车采用相同程序；若制动能力不同，不能再直接相消。

\textbf{后车制动更强时，最小间距未必在停车后。} 令前车制动减速度为
\(b_F\)，后车为 \(b_R>b_F\)。等待时间 \(\tau\) 内，后车仍以 \(v\)
匀速。定义已经损失的间距 \[
L(t)=\int_0^t(v_R-v_F)\,dt.
\] 安全条件应写成
\(D-\max_t L(t)>d_*\)，而不能只检查两车停车位置。两车都未停的阶段， \[
v_R-v_F=b_R\tau-(b_R-b_F)t.
\] 其零点为 \[
t_*=\frac{b_R\tau}{b_R-b_F}.
\] 若
\(t_*<v/b_F\)，则后车在前车停下前追平速度，随后速度小于前车，间距开始增加。由此得到反应时间的分界
\[
\tau_c=v\left(\frac1{b_F}-\frac1{b_R}\right).
\] 对于 \(0\le\tau<\tau_c\)，先积等待段，再积共同制动段： \[
L(\tau)=\frac12b_F\tau^2,
\] \[
L_{\max}
=\frac12b_F\tau^2+
\frac{(b_F\tau)^2}{2(b_R-b_F)}
=\boxed{\frac{b_Fb_R\tau^2}{2(b_R-b_F)}}.
\]
这里第二项是速度差图上的三角形面积。因其对应的时间仍在两车共同制动阶段，不能把前车停车后的直角三角形错误加进来。

对于
\(\tau\ge\tau_c\)，后车不会在前车停车前降到相同速度。直到后车也停下，\(L\)
单调增加，最大值为 \[
L_{\max}
=v\tau+\frac{v^2}{2b_R}-\frac{v^2}{2b_F}.
\] 两支在 \(\tau_c\) 连续： \[
L_{\max}(\tau_c)=
\frac{v^2(b_R-b_F)}{2b_Fb_R}.
\] 一阶导数也连续，两边均为
\(v\)。这个检查可以发现漏掉等待距离或写反停止距离的错误。

例如
\(v=10\,\mathrm{m/s}\)、\(b_F=1\,\mathrm{m/s^2}\)、\(b_R=2\,\mathrm{m/s^2}\)、\(\tau=2\,\mathrm s\)，有
\(\tau_c=5\,\mathrm s\)， \[
L_{\max}=4\,\mathrm m,\qquad t_*=4\,\mathrm s.
\] 停车后的损失却为 \[
L_{\rm final}=20+25-50=-5\,\mathrm m.
\] 停车后间距比最初更大，过程中仍曾缩小
\(4\,\mathrm m\)。这说明用停车位置代替全程最小间距可能给出相反的判断。图1-2的等制动能力模型没有这一转折，不能把它的单调性直接搬到强制动后车上。

\textbf{速度上限变化时，启动门槛如何决定延迟。} 令门槛增量为
\(s=d_s-d_0\)。前车的加速距离为 \(s_a=v^2/(2a)\)。统一启动延迟为 \[
\Delta=
\begin{cases}
\sqrt{2s/a},&s\le s_a,\\
s/v+v/(2a),&s>s_a.
\end{cases}
\] 在 \(s=s_a\)，两式均为 \(v/a\)。两车完成启动之后，间距恒为
\(d_0+v\Delta\)。但启动时门槛为
\(d_0+s\)，这两个量通常不同；达到限速后的间距还包含后车加速过程中形成的距离增量。

\needspace{14\baselineskip}

在第一支，\(\partial \Delta/\partial s=1/\sqrt{2as}\)，小门槛时等待时间对门槛比较敏感；在第二支，导数为
\(1/v\)。这只是模型对参数的敏感性，不能据此把司机反应时间与车辆启动延迟混用。启动延迟由距离门槛产生，制动延迟由反应程序产生，两者只有在题设专门规定时才相等。

\textbf{带有限车长的队列。} 若每辆车长度为
\(\ell_c\)，匀速时前车尾到后车头为 \(d_\infty\)，则相邻车头的时间间隔为
\[
t_{\rm head}=\frac{\ell_c+d_\infty}{v}.
\] 第 \(N\) 辆车在第 \((N-1)\Delta\)
时启动，但这不是匀速通过固定截面的车头间隔。比较这两个时间前，应先注明测量的是启动事件还是经过事件。若车长不同，题1中相邻尾---头间距的求法仍成立，车头间隔则需逐辆加入相应车长。

\subsection{六、关键洞察}\label{ux516dux5173ux952eux6d1eux5bdf}

车距变化只由相对速度的面积决定。启动门槛决定时间差；限速决定这段时间差最终对应的空间差。汽车长度、各车的绝对位置和同样的加速位移都在相减中消去。

制动题中``始终大于''不能未经检验就改成``最后大于''。本题之所以可以这样做，是因为前车先进入同样的减速过程，始终有
\(v_F\le v_R\)。间距单调性给出了完整的安全条件。

\subsection{七、易错点}\label{ux4e03ux6613ux9519ux70b9}

把 \(7\,\mathrm m\) 当成最终间距；把 \(7\) 代入启动位移而忘记减去初始的
\(2\)；后车启动后仍把它视为静止；用车头---车头距离替换题目定义；将严格上限写成允许等号；未证明最小距离发生在停车之后。

\subsection{八、课堂讲解建议}\label{ux516bux8bfeux5802ux8bb2ux89e3ux5efaux8bae}

先让学生在图1-1上圈出``前车尾''和``后车头''，再用两色画速度曲线。第一遍解题只算两块面积；第二遍写分段函数，检查面积解法是否遗漏阶段。

比较 \(\tau=0\)、\(\tau=\Delta/2\)、\(\tau=\Delta\)
三种情形。相应最终距离分别为 \(d_\infty\)、\(d_0+v\Delta/2\) 和
\(d_0\)。这种参数检查能直接辨认不等式的方向，也能说明等待越久，最终距离越小。

制动开始时，应另设时间零点。沿用启动时的时间又在式中直接写
\(v-bt\)，会混淆两段运动的起始时刻。

\textbf{启动事件沿车队怎样传播。} 设每车长度为
\(\ell\)，初始相邻车的位置标记相差 \(\ell+d_0\)。第 \(j\) 辆车在
\(t_j=j\Delta\) 启动，它在启动前尚未移动，因此启动事件坐标为
\(x_j=-j(\ell+d_0)\)。事件在位置---时间图上的斜率给 \[
v_{\rm wave}=-\frac{\ell+d_0}{\Delta}.
\]
它向队尾传播，与启动后汽车向前行驶的速度不同。若车长未给，原题仍可求车距和启动延迟，却不能给这个传播速度的数值。该变式能说明为何车长在原题间距计算中消去，在涉及事件传播的任务中却必须保留。

\subsection{九、变式训练}\label{ux4e5dux53d8ux5f0fux8badux7ec3}

\textbf{变式一：两车制动能力不同。} 前车制动加速度为 \(b_F\)，后车延迟
\(\tau\)，随后以 \(b_R\le b_F\) 制动，初速度均为 \(v\)。求初始间距 \(D\)
必须满足的条件，使停车后间距大于 \$d\_* \$。

后车制动不更强，且更晚开始，故整个过程中其速度不低于前车。最终间距就是最小间距：
\[
D_{\min}=D-v\tau-\frac{v^2}{2b_R}+\frac{v^2}{2b_F}.
\] 因此 \[
\boxed{D>d_*+v\tau+\frac{v^2}{2b_R}-\frac{v^2}{2b_F}}.
\] 若 \(b_R>b_F\)，相对速度可能中途变号，最小距离可能早于停车，必须另求
\(v_F=v_R\) 的时刻；仅看停车距离不再充分。

\textbf{变式二：启动门槛超过加速距离。} 若
\(d_s-d_0>v^2/(2a)\)，前车已进入匀速阶段时后车才启动。此时 \[
d_s-d_0=v\Delta-\frac{v^2}{2a},
\] 所以 \[
\boxed{\Delta=\frac{d_s-d_0}{v}+\frac{v}{2a}}.
\] 最终间距仍是
\(d_0+v\Delta\)，但不能再用平方根公式求延迟。先判断门槛所在阶段，是分段运动中比代数运算更早的一步。

\textbf{课堂核验：把``停得更近''与``途中追得更近''分开。}
可先给上述强制动例，要求学生画
\(v_R-v_F\)，只标等待结束、速度追平、前车停车、后车停车四个事件。相对速度的正负决定间距增减；曲线下面积决定增减多少。仅列两辆车的停车距离不能证明途中安全。

若用计算表检验，时间步应穿过 \(t=\tau,t_*,v/b_F,\tau+v/b_R\)，并同时记录
\(d(t)\)。正确结果在 \(t_*\)
出现最小值，而不是在最后一行出现最小值。数值表用于发现分段次序错误，解析三角形面积仍是答案的依据。

\textbf{含前车初始速度差的迁移题。} 在后车开始制动的时刻，设前车速度为
\(u_F\)、后车为 \(u_R\)，两车均未停止且
\(u_R>u_F\)、\(b_R>b_F\)。相对速度为 \[
w(t)=u_R-u_F-(b_R-b_F)t.
\] 若速度追平发生在任一辆车停车之前，则此阶段损失的间距为 \[
\boxed{\Delta L=\frac{(u_R-u_F)^2}{2(b_R-b_F)}}.
\] 条件应逐项检查： \[
\frac{u_R-u_F}{b_R-b_F}
<\min\left(\frac{u_F}{b_F},\frac{u_R}{b_R}\right).
\]
条件不成立时需要换成停车后的速度表达式。这一小题可以检查学生是否真正理解了分段模型，而不是记住某一个
\(v\tau\) 的结论。

\needspace{14\baselineskip}

\textbf{完整分段速度差的检验。}
等制动能力情形，后车反应段速度差随时间从零增大，后车开始刹车后速度差恒定，前车停车后再下降到零。其全程积分仍为
\(v\tau\)。当后车制动更强时，中段速度差斜率转负，最小间距会提前出现。板书可以把这两个图并排，不先计算停车距离。

对于
\(\tau\ge v/b_F\)，前车已在后车刹车前停止。等待段应再分为前车制动和前车静止两段。直接用``共同制动''公式跨过前车停车时刻会出现负速度；但通过总位移公式求最终损失依然正确。独立数值核验应使用
\(\max(v-bt,0)\)，而非无限延伸的直线。

本题严格安全要求间距大于阈值。如果在临界值正好等于阈值，可称边界，不能同时把等号计为``严格大于''。末态与过程中最小值应使用同一口径。

\newpage

\section{第 2 题}\label{ux7b2c-2-ux9898}

用德布罗意波和驻波条件研究受限制的电子，电子质量为 \(m\)，普朗克常量为
\(h\)。

（1）电子限制在间距为 \(L\) 的两堵理想硬墙之间，接触点为驻波节点。

\begin{enumerate}
\def\labelenumi{\alph{enumi}.}
\tightlist
\item
  求量子化后的速度大小和能量；b.
  求电子推一面墙的平均力；c.~求位移---动量相空间中的面积。
\end{enumerate}

（2）电子限制在半径为 \(R\) 的圆形硬墙内。分别求沿正三角形、正方形和正
\(\ell\) 边形周期轨道运动时的能量 \(E_{3n},E_{4n},E_{\ell n}\)。

\includegraphics[width=0.88\linewidth,height=\textheight,keepaspectratio,alt={题2　硬墙驻波与圆内周期多边形}]{../figures/question-02.pdf}

\subsection{一、题型定位}\label{ux4e00ux9898ux578bux5b9aux4f4d-1}

前半题是一维无限深势阱，可由驻波条件得到精确能级；后半题用周期轨道量子化作模型计算。两部分共享德布罗意关系
\(p=h/\lambda\)，但量子化路径和适用范围不同。

\subsection{二、知识点}\label{ux4e8cux77e5ux8bc6ux70b9-1}

驻波节点；动量大小与有向动量的区别；碰撞冲量的周期平均；能级对边界位置的导数；作用量
\(\oint p\,dx\)；闭合路径的总相位；半经典模型与完整本征值问题。

\subsection{三、物理图像}\label{ux4e09ux7269ux7406ux56feux50cf-1}

一维墙间的第 \(n\) 个驻波含 \(n\)
个半波长。经典碰撞图像中，电子在两墙之间往返，动量在 \(+p_n\) 与
\(-p_n\) 之间变号。图2-2用矩形表示作用量围成的面积。

圆内的正多边形每一条边都连接相邻反射点，入射角与反射角相等。多边形总周长决定沿一整条周期轨道累积的传播相位；它不是圆盘内所有可能波函数的分类。

\subsection{四、解题入口}\label{ux56dbux89e3ux9898ux5165ux53e3-1}

一维先写 \(L=n\lambda_n/2\)，再求
\(p_n,v_n,E_n\)。推墙力可从冲量频率计算，并用 \(-dE_n/dL\)
复核。圆内轨道先算周长，再应用题目采用的总相位条件。

\subsection{五、解析}\label{ux4e94ux89e3ux6790-1}

\textbf{一维驻波的允许波长。} 两端为节点，一段长度 \(L\)
内放入整数个半波长： \[
L=n\frac{\lambda_n}{2},\qquad n=1,2,\ldots .
\] 因此 \[
\lambda_n=\frac{2L}{n},
\qquad
\boxed{p_n=\frac{nh}{2L}}.
\] 这里 \(p_n\)
表示动量大小。驻波本征态的平均有向动量为零，不代表它的动能为零；其动量平方的期望为
\(p_n^2\)。

在非相对论范围 \(p_n\ll mc\)， \[
\boxed{v_n=\frac{nh}{2mL}},
\qquad
\boxed{E_n=\frac{n^2h^2}{8mL^2}}.
\] 波函数可以写成 \[
\psi_n(x)=\sqrt{\frac2L}\sin\frac{n\pi x}{L},
\qquad0<x<L.
\] 它在 \(x=0,L\) 均为零，归一化积分为一： \[
\int_0^L|\psi_n|^2dx
=\frac2L\int_0^L\sin^2\frac{n\pi x}{L}\,dx=1.
\] 直接代入动能算符 \[
-\frac{\hbar^2}{2m}\frac{d^2}{dx^2}
\] 同样得到 \(E_n=n^2\pi^2\hbar^2/(2mL^2)\)，与驻波结果一致，其中
\(\hbar=h/(2\pi)\)。

\textbf{一面墙上的平均力。}
在经典往返图像中，电子碰到同一面墙的时间间隔为 \[
\Delta t=\frac{2L}{v_n}.
\] 一次弹性反射时，电子动量改变
\(2p_n\)，墙获得大小相同、方向相反的冲量。因此一面墙的平均力为 \[
\begin{aligned}
F_n
&=\frac{2p_n}{2L/v_n}
=\frac{p_nv_n}{L}
=\frac{p_n^2}{mL}\\
&=\boxed{\frac{n^2h^2}{4mL^3}}
=\frac{2E_n}{L}.
\end{aligned}
\] 两墙的受力大小相同、方向相反。分母应是回到同一面墙的时间
\(2L/v_n\)，不是走到另一面墙的 \(L/v_n\)；后者会多算一倍。

\textbf{从能级变化复核推墙力。} 将右墙缓慢向右移动
\(dL\)，且电子保持在同一能级 \(n\)。电子对墙做功
\(F_n dL\)，其能量降低： \[
dE_n=-F_n\,dL.
\] 因而 \[
F_n=-\frac{d}{dL}\left(\frac{n^2h^2}{8mL^2}\right)
=\frac{n^2h^2}{4mL^3}.
\]
两种推导分别使用碰撞频率和能级对长度的响应，结果一致。量子本征态并不需要电子沿某条确定轨迹不断碰墙；碰撞图像只是帮助理解同一平均力的机械模型。

\begin{figure}[H]
\centering
\includegraphics[width=0.9\linewidth,height=\textheight,keepaspectratio,alt={图2-2　作用量面积与正多边形周长}]{../figures/solution-02.pdf}
\caption*{图2-2　作用量面积与正多边形周长}
\end{figure}

\textbf{相空间的面积。} 以 \(x\) 为横轴、\(p_x\) 为纵轴，经典能量不超过
\(E_n\) 的允许区域为 \[
0\le x\le L,
\qquad -p_n\le p_x\le p_n.
\] 其面积 \[
\boxed{\mathcal A_n=L(2p_n)=nh}.
\] 同一结果可由沿一次往返轨道的作用量求得： \[
\oint p_x\,dx
=\int_0^L p_n\,dx+\int_L^0(-p_n)\,dx
=2Lp_n=nh.
\]
这里的面积是两条有向动量支路围出的作用量面积，或能量上界确定的相空间区域面积。若只讨论严格固定能量的经典轨迹，轨迹本身由两条线段组成，二维面积为零；不能把这两种``范围''混为一谈。量子态也不是在矩形中每一点具有同样概率。

\textbf{圆内正多边形的几何长度。} 正 \(\ell\) 边形内接于半径 \(R\)
的圆，每边对应圆心角 \(2\pi/\ell\)。取半条边构成的直角三角形： \[
\frac{s_\ell}{2}=R\sin\frac\pi\ell.
\] 所以边长与周长为 \[
s_\ell=2R\sin\frac\pi\ell,
\qquad
P_\ell=2\ell R\sin\frac\pi\ell.
\] 题目采用的闭合轨道条件是完整周长容纳整数个德布罗意波长： \[
P_\ell=n\lambda_{\ell n}.
\] 因此 \[
p_{\ell n}=\frac{nh}{P_\ell},
\] \[
\boxed{E_{\ell n}=
\frac{n^2h^2}{8m\ell^2R^2\sin^2(\pi/\ell)}}.
\] 原题未给出圆内周期轨道的量子化条件，这里明确采用 \(P_\ell=n\lambda\)
的补充模型。所得数值系数并非由原题的几何条件单独确定。这是忽略反射附加相位、横向模式以及轨道族干涉的周期轨道模型结果。若把每个反射点都当成一维驻波节点，就会另写
\(s_\ell=n\lambda/2\)；它对应不同的额外假设，不能在不说明的情况下替换总路径条件。

\textbf{三角形与正方形。} 对 \(\ell=3\)， \[
P_3=6R\sin60^\circ=3\sqrt3R,
\] 故 \[
\boxed{E_{3n}=\frac{n^2h^2}{54mR^2}}.
\] 对 \(\ell=4\)， \[
P_4=8R\sin45^\circ=4\sqrt2R,
\] 故 \[
\boxed{E_{4n}=\frac{n^2h^2}{64mR^2}}.
\] 在同一 \(n\) 下，正方形周长更大，所以动量和能量更小。这与能量按
\(P_\ell^{-2}\) 缩放相符。

\textbf{大边数极限。} 当 \(\ell\to\infty\)， \[
\ell\sin\frac\pi\ell\to\pi,
\qquad P_\ell\to2\pi R.
\] 于是 \[
E_{\ell n}\longrightarrow
\frac{n^2h^2}{8\pi^2mR^2}
=\frac{n^2\hbar^2}{2mR^2}.
\]
这与仅沿圆周运动的环上自由粒子能级形式相同，但不把圆盘变成一维圆环。圆盘仍允许径向运动，完整边界条件需要同时求解径向与角向波函数。

\textbf{完整圆盘问题为什么多一个量子数。}
圆内自由电子的二维定态方程在极坐标中可分离为 \[
\psi(r,\vartheta)=F(r)e^{ij\vartheta},\qquad j\in\mathbb Z.
\] 正则径向解为 \(F(r)=J_{|j|}(kr)\)。硬墙要求 \(F(R)=0\)，即 \[
kR=z_{|j|,s},
\] 其中 \(z_{|j|,s}\) 是第 \(s\) 个正贝塞尔零点。能级为 \[
E_{j,s}=\frac{\hbar^2z_{|j|,s}^2}{2mR^2}.
\]
这个式子用于说明模型边界：周期三角形或正方形不是完整圆盘本征态的精确能级标签。原题给出的轨道形式只能在所采用的简化量子化方案内使用。

\textbf{从波函数检查``推墙力''并非粒子固定站在墙上。}
一维无限深势阱的归一化定态为 \[
\psi_n(x)=\sqrt{\frac2L}\sin\frac{n\pi x}{L}.
\] 其概率密度在墙处为零，但导数不为零： \[
\psi_n'(L)=(-1)^n\sqrt{\frac2L}\frac{n\pi}{L}.
\] 由硬壁上的量子动量通量，右壁所受平均力为 \[
F_n=\frac{\hbar^2}{2m}|\psi_n'(L)|^2
=\frac{n^2\pi^2\hbar^2}{mL^3}.
\] 代入 \(h=2\pi\hbar\)，它与 \(2E_n/L\)
一致。墙处的粒子概率为零，不代表作用力为零；硬壁改变波函数的空间梯度，动量通量由这个梯度决定。

\textbf{缓慢移动与突然移动的物理过程不同。} 若右墙缓慢从 \(L_1\) 移到
\(L_2\)，且满足绝热条件，量子数保持不变： \[
\frac{E_f}{E_i}=\left(\frac{L_1}{L_2}\right)^2,\qquad
W_{\rm by}=E_i-E_f.
\] 外界压缩时
\(L_2<L_1\)，\(W_{\rm by}<0\)，即外界对电子做正功。对每个瞬时长度， \[
F(L)=\frac{n^2h^2}{4mL^3}.
\] 积分给 \[
\int_{L_1}^{L_2}F(L)\,dL
=\frac{n^2h^2}{8m}
\left(\frac1{L_1^2}-\frac1{L_2^2}\right),
\] 与能量变化相符。

若墙突然由 \(L\) 扩到 \(2L\)，改变瞬间波函数尚未来得及变化。以新定态
\(\phi_j(x)\) 展开， \[
c_j=\int_0^L
\sqrt{\frac1L}\sin\frac{j\pi x}{2L}
\sqrt{\frac2L}\sin\frac{n\pi x}{L}\,dx.
\] 对 \(j\ne2n\)， \[
c_j=
\frac{4\sqrt2\,n(-1)^{n+1}\sin(j\pi/2)}
{\pi(4n^2-j^2)},
\] 而 \(c_{2n}=1/\sqrt2\)。所有偶数 \(j\ne2n\)
的系数为零，奇数能级通常有非零占据。扩张后不是一个单独的第 \(n\)
能级，所以不能使用缓慢扩张的 \(E_i/4\)。

突然扩张的理想硬壁模型中，原波函数在 \(x>L\) 延拓为零，其动能二次型仍为
\[
\frac{\hbar^2}{2m}\int_0^{2L}|\psi'(x)|^2dx=E_i.
\] 于是
\(\sum_j|c_j|^2E_j(2L)=E_i\)。平均推墙力在后续时间会随态的相干演化变化；若只取长时间去相干平均，
\[
\overline F=\sum_j|c_j|^2\frac{2E_j(2L)}{2L}
=\frac{E_i}{L}.
\] 它是原定态平均力的一半，而缓慢扩张到 \(2L\) 的推墙力为原值的
\(1/8\)。同样的最终几何尺寸，因过程不同，量子态和受力也不同。

\needspace{14\baselineskip}

\textbf{多个互不作用电子占据不同能级。}
若允许两种自旋态，每个空间能级最多容纳两个电子。对 \(N_e=2K\)
个电子的零温基态， \[
U=2\sum_{n=1}^K E_n
=\frac{h^2}{4mL^2}\frac{K(K+1)(2K+1)}6.
\] 总推墙力为 \[
F_{\rm tot}=-\frac{\partial U}{\partial L}
=\frac{2U}{L}.
\] 若 \(N_e=2K+1\)，再加入一个第 \(K+1\) 级电子： \[
U_{\rm odd}=U_{\rm even}+E_{K+1},\qquad
F_{\rm odd}=2U_{\rm odd}/L.
\] 不能把所有电子都放在
\(n=1\)，否则漏掉费米统计。这个模型忽略电子之间的库仑相互作用，适合检查量子态计数，并不是实际电子气全部相互作用的答案。

\textbf{相空间面积如何用于态数估计。} 图2-2中的矩形面积为
\(2Lp\)。一维每个自旋态的半经典态数约为 \[
\mathcal N(E)\simeq\frac{2L\sqrt{2mE}}h.
\] 在硬壁精确谱中，空间态数是 \[
\mathcal N(E)=
\left\lfloor\frac{2L\sqrt{2mE}}h\right\rfloor.
\] 面积除以 \(h\)
给出连续近似；向下取整给出实际整数态数。二维圆盘则需四维相空间体积，不能把某条正多边形轨道的周长量子化当成全部圆盘态数。

\subsection{六、关键洞察}\label{ux516dux5173ux952eux6d1eux5bdf-1}

量子化首先是边界或相位条件，再是把允许动量代入动能公式。一维两节点产生半波长条件；闭合周期路径使用整波长条件。只背
\(E\propto n^2\)，容易丢掉两者不同的因子。

推墙力的两种算法也揭示了压力的含义：碰撞给出机械冲量，能级导数给出量子系统对边界变化的响应。在一维硬墙模型中，它们可以相互校验。

\subsection{七、易错点}\label{ux4e03ux6613ux9519ux70b9-1}

把驻波态的平均动量零误解为动能零；同一面墙的碰撞间隔少写一倍；把相空间矩形当成量子概率均匀分布；圆内边长误用圆弧长；把正多边形轨道能级当作圆盘精确谱；不说明反射相位假设就更改量子化条件。

\subsection{八、课堂讲解建议}\label{ux516bux8bfeux5802ux8bb2ux89e3ux5efaux8bae-1}

在图2-1上数半波，在图2-2上沿正、负动量支路各走一遍。往返作用量中第二段的
\(p\) 和 \(dx\) 都为负，乘积仍为正，这一步能同时解释面积与因子2。

讲圆内轨道时，先只做几何计算，再将 \(P_\ell\)
写入相位条件。让学生分别用三角形、正方形和大边数极限检查通式。最后指出``看起来像经典轨道''与``完整波函数满足二维边界''是不同层次的问题。

\textbf{突然扩阱后的平均能量与方差。}
旧阱本征波函数在扩大的空间中仍连续，但旧墙位置处一阶导数跳变。对新阱本征态展开，大量子数系数具有
\(c_j=O(j^{-2})\) 的尾部（某些对称支系数为零）。新能量
\(E_j\propto j^2\)，所以 \[
\sum_j E_j|c_j|^2\sim\sum_jj^{-2}
\] 收敛，平均能量有限；而 \[
\sum_j E_j^2|c_j|^2
\]
通常不收敛。理想瞬时移动硬墙可以产生无限能量方差，即使平均能量守恒。物理墙的有限移动时间和有限势垒会改变高能尾部；这说明仅检查平均能量并不能确认全部高阶矩都良好。

\subsection{九、变式训练}\label{ux4e5dux53d8ux5f0fux8badux7ec3-1}

\textbf{变式一：缓慢压缩硬墙。} 墙间距离由 \(L\) 缓慢变为
\(L/2\)，电子一直保持第 \(n\) 个能级，求能量增量和外界做功。

由 \(E_n\propto L^{-2}\)，终态能量为 \(4E_n(L)\)。外界功 \[
\boxed{W_{\rm on}=3E_n(L)}.
\] 也可积分推墙力： \[
W_{\rm on}=-\int_L^{L/2}F_n(\ell)\,d\ell
=\frac{n^2h^2}{8m}\left[\frac1{(L/2)^2}-\frac1{L^2}\right].
\]
``缓慢''用于保持能级；突然移动墙时，初态通常投影到多个终态能级，不能直接沿用这一结果。

\textbf{变式二：多边形周长的误差。} 对很大的
\(\ell\)，求能量相对圆周极限的首项修正。 \[
\ell\sin\frac\pi\ell
=\pi\left[1-\frac{\pi^2}{6\ell^2}+O(\ell^{-4})\right].
\] 周长平方的倒数给出 \[
\boxed{E_{\ell n}=
\frac{n^2\hbar^2}{2mR^2}
\left[1+\frac{\pi^2}{3\ell^2}+O(\ell^{-4})\right]}.
\]
有限边数的折线比圆周短，所以能量比极限值高；修正项正号符合这一几何判断。

\textbf{评讲突然扩张时，应把能级、占据与总能量分开。} 墙变为 \(2L\)
后，新的能级全部下降，但电子并没有自动进入新基态。新旧本征函数的重叠决定占据，多个新能级共同维持原来的瞬时平均动能。若学生说``长度翻倍，任何状态能量都变成四分之一''，可以用
\(c_{2n}=1/\sqrt2\) 立即检验：仅这一项已贡献
\(E_i/2\)，其余奇数态还贡献能量。

\textbf{对圆形轨道量子化的审题要求。}
正多边形只描述一族经典周期轨道。若加入反射相位
\(\chi\)，传播闭合条件成为 \[
\frac{pP_\ell}{\hbar}+\ell\chi=2\pi n.
\] 所以 \[
p=\frac h{P_\ell}\left(n-\frac{\ell\chi}{2\pi}\right).
\] 原稿采用的 \(pP_\ell=nh\)
对应忽略或已吸收反射相位。课堂上可把这个条件作为明确补充模型，不能告诉学生它从圆内几何自动推出。若指定硬壁的完整波函数边界，应用贝塞尔零点；两类问题的状态编号和相位条件不同。

\textbf{一个可核对的边界压力变式。} 对固定概率占据 \(w_n\)
的混合态，\(\sum w_n=1\)，缓慢移动壁而占据不变时 \[
U=\sum_nw_nE_n,\qquad F=2U/L.
\] 若连接热库保持温度，占据 \(w_n(L,T)\)
随长度变化，应从自由能求力，不能把 \(-\partial U/\partial L\)
直接当成等温力。固定占据与固定温度的导数条件是这一变式的评分点。

\needspace{14\baselineskip}

\textbf{突然扩箱中的能量口径。} 壁突然从 \(L\) 移到 \(2L\)
时，原波函数作为新空间中的函数，旧区外为零；瞬间没有时间完成绝热能级跟随。它在新本征态中的展开并非只有第
\(2n\) 项，尽管该项能量等于原
\(E_n\)。奇数项也有非零权重，它们的总体使能量期望仍为原值。

系数平方的归一化 \(\sum_j|c_j|^2=1\) 与能量和 \(\sum_j|c_j|^2E_j'=E_n\)
是不同检查。截断求和的能量尾部比归一化尾部下降慢，数值验证时不能给两者同一个截断误差尺度。

圆内闭合多边形只代表一条经典轨道的相位条件；完整硬壁圆盘本征函数涉及径向与角向模式。这一模型层级应在题目答案中直接注明，避免学生把几何轨道长度的量子化误称圆盘全部谱。

\newpage

\section{第 3 题}\label{ux7b2c-3-ux9898}

均匀膨胀满足 \(v(t)=H(t)r(t)\)。以当前距离 \(r_0\) 定义尺度因子
\(a(t)=r(t)/r_0\)，当前时刻为 \(t_0\)，有 \(a(t_0)=1\)。把大爆炸时刻记为
\(t=0\)，取极限 \(a(t)\to0\)。

（1）用 \(a,\dot a\) 表示 \(H(t)\)。

（2）只考虑均匀质量密度 \(\rho\) 的万有引力，求 \(\ddot a/a\)。

（3）在膨胀与再坍缩的临界、总机械能为零的模型中，求 \(\dot a^2\)。

（4）若共动球内总质量守恒，用当前密度 \(\rho_0\) 与 \(a(t)\) 表示
\(\rho(t)\)。

（5）若 \(a(t)=(t/t_0)^n\)，求 \(n\)。

（6）给定
\(H_0=72\,\mathrm{km\,s^{-1}\,Mpc^{-1}}\)，\(1\,\mathrm{Mpc}=10^6\,\mathrm{pc}\)，\(1\,\mathrm{pc}=3.09\times10^{16}\,\mathrm m\)，求本模型的当前年龄。

（7）求当前密度 \(\rho_0\)。

\includegraphics[width=0.84\linewidth,height=\textheight,keepaspectratio,alt={题3　共动球与尺度因子}]{../figures/question-03.pdf}

\subsection{一、题型定位}\label{ux4e00ux9898ux578bux5b9aux4f4d-2}

这是均匀无压物质的牛顿宇宙学模型。哈勃关系是运动学定义，球壳定理提供动力学，质量守恒给出密度随尺度的变化。总能量为零后，尺度因子才被确定为时间的
\(2/3\) 次幂。

\subsection{二、知识点}\label{ux4e8cux77e5ux8bc6ux70b9-2}

共动距离；球壳定理；零机械能条件；质量密度随体积变化；幂律函数；单位换算；临界密度。模型不包含辐射压力、宇宙常量和暗能量。

\subsection{三、物理图像}\label{ux4e09ux7269ux7406ux56feux50cf-2}

图3-1跟踪同一批物质围成的球，其边界半径是
\(r=a r_0\)。球内质量不随膨胀改变，体积却按 \(a^3\)
增长，所以密度下降。图3-2中 \(a\)
持续增大而曲线向下弯，表示膨胀仍在继续，但速度逐渐减小。

\subsection{四、解题入口}\label{ux56dbux89e3ux9898ux5165ux53e3-2}

对球边界上的单位质量粒子写引力加速度与机械能。先保留
\(r(t)\)，完成球对称计算后再用 \(r=a r_0\) 消去任意选取的
\(r_0\)，可避免把尺度因子误当成有长度单位的半径。

\subsection{五、解析}\label{ux4e94ux89e3ux6790-2}

\textbf{哈勃变量。} 因为 \(r_0\) 是固定的共动标签， \[
r=a r_0,\qquad \dot r=\dot a r_0.
\] 代入 \(\dot r=Hr\)，在 \(a>0\) 时得到 \[
\boxed{H(t)=\frac{\dot a}{a}}.
\] 它的量纲为逆时间。\(H\)
可以随时间改变；``速度与距离成正比''并不要求比例系数在所有时刻相同。

\textbf{引力加速度。} 均匀球内部总质量 \[
M(r)=\frac{4\pi}{3}\rho r^3.
\]
球壳定理表明，边界粒子的径向引力来自球内质量，球外同心壳不贡献净力。因此
\[
\ddot r=-\frac{GM(r)}{r^2}
=-\frac{4\pi G}{3}\rho r.
\] 用 \(r=a r_0\) 消去 \(r_0\)： \[
\boxed{\frac{\ddot a}{a}=-\frac{4\pi G}{3}\rho}.
\]
负号表示引力使外向膨胀减速，不表示尺度因子必须立刻减小。是否最终回缩，还要由初速度和总能量判定。

\textbf{临界零能量。} 共动球内的质量 \(M\)
固定，单位质量边界粒子的机械能为 \[
\mathcal E=\frac12\dot r^2-\frac{GM}{r}.
\] 在临界情形 \(\mathcal E=0\)： \[
\dot r^2=\frac{2GM}{r}.
\] 代入 \(M=4\pi\rho r^3/3\)，得到 \[
\dot r^2=\frac{8\pi G}{3}\rho r^2.
\] 再以 \(r_0^2\) 除两边： \[
\boxed{\dot a^2=\frac{8\pi G}{3}\rho a^2},
\qquad
\boxed{H^2=\frac{8\pi G}{3}\rho}.
\] 这个等式只适用于题目指定的零能量物质模型。若总能量非零，还会有额外的
\(a^{-2}\) 项。

\begin{figure}[H]
\centering
\includegraphics[width=0.837\linewidth,height=\textheight,keepaspectratio,alt={图3-2　物质模型的尺度、哈勃变量与密度}]{../figures/solution-03.pdf}
\caption*{图3-2　物质模型的尺度、哈勃变量与密度}
\end{figure}

\textbf{质量守恒。} 同一共动球在当前时刻的质量为 \[
M=\frac{4\pi}{3}\rho_0r_0^3.
\] 在时刻 \(t\)， \[
M=\frac{4\pi}{3}\rho(t)[a(t)r_0]^3.
\] 两式相等，给出 \[
\boxed{\rho(t)=\rho_0a^{-3}(t)}.
\] 代回零能量式： \[
\dot a^2=\frac{8\pi G\rho_0}{3}\frac1a.
\] 记 \[
K=\frac{8\pi G\rho_0}{3}>0.
\] 膨胀支路取 \(\dot a=\sqrt K a^{-1/2}\)。变量分离： \[
a^{1/2}da=\sqrt K\,dt.
\] 从大爆炸极限积分到时刻 \(t\)： \[
\frac23a^{3/2}=\sqrt K\,t.
\] 因此 \[
a(t)=\left(\frac32\sqrt K\,t\right)^{2/3}.
\] 当前归一化 \(a(t_0)=1\) 给出 \[
\sqrt K=\frac{2}{3t_0},
\] 所以 \[
\boxed{a(t)=\left(\frac{t}{t_0}\right)^{2/3}},
\qquad\boxed{n=\frac23}.
\] 这条积分路线同时确定指数和归一化常数，比只比较幂次更完整。

\textbf{用幂次匹配复算。} 设 \(a=(t/t_0)^n\)，则 \[
\dot a^2\propto t^{2n-2},
\qquad a^{-1}\propto t^{-n}.
\] 零能量式要求 \[
2n-2=-n,
\] 得到 \(n=2/3\)。再求 \[
\boxed{H(t)=\frac{2}{3t}},
\qquad
\boxed{\rho(t)=\frac{1}{6\pi Gt^2}}.
\] 第二式也可以从 \(\rho_0a^{-3}\) 得到。将 \(a\propto t^{2/3}\)
代入加速度式，有 \[
\frac{\ddot a}{a}=-\frac{2}{9t^2}
=-\frac{4\pi G}{3}\frac{1}{6\pi Gt^2},
\] 说明动力学与能量计算完全一致。

\textbf{当前年龄。} 原题给出的长度换算为 \[
1\,\mathrm{Mpc}=3.09\times10^{22}\,\mathrm m.
\] 因此 \[
H_0=\frac{72\times10^3}{3.09\times10^{22}}\,\mathrm{s^{-1}}
\simeq2.33010\times10^{-18}\,\mathrm{s^{-1}}.
\] 由 \(H_0=2/(3t_0)\)， \[
\boxed{t_0=\frac{2}{3H_0}\simeq2.86111\times10^{17}\,\mathrm s}.
\] 采用 \(1\,\mathrm{yr}=365.25\times86400\,\mathrm s\)，得到 \[
\boxed{t_0\simeq9.07\times10^9\,\mathrm{yr}}.
\] 逆哈勃量 \(H_0^{-1}\simeq13.60\times10^9\,\mathrm{yr}\)
是时间尺度，不能直接当成本题的年龄；物质零能量模型还有 \(2/3\)
因子。这里求的是题设模型的年龄，不是用现代多组分宇宙模型拟合出的观测年龄。

\textbf{当前密度。} 在当前时刻，零能量式为 \[
\boxed{\rho_0=\frac{3H_0^2}{8\pi G}}.
\] 以 \(G=6.67\times10^{-11}\,\mathrm{m^3\,kg^{-1}\,s^{-2}}\)
作为数值计算常量，得 \[
\boxed{\rho_0\simeq9.72\times10^{-27}\,\mathrm{kg/m^3}}.
\] 原卷没有另给 \(G\)
的数值，因此符号式是基本答案；数值需要注明所用常量。量纲检查为 \[
\frac{H_0^2}{G}
\sim\frac{\mathrm{s^{-2}}}{\mathrm{m^3\,kg^{-1}\,s^{-2}}}
=\mathrm{kg/m^3}.
\]

\textbf{大爆炸端点的含义。} 在 \(t\downarrow0\) 时，\(a\downarrow0\)，而
\(H,\rho\) 发散。因此 \(a(0)=0\) 是解的奇异端点，不能把 \(t=0\)
当成普通的有限密度初始条件，直接代入 \(H=\dot a/a\)。所有微分等式应在
\(t>0\) 使用，再讨论端点极限。

\textbf{临界与非临界的区别。} 对非零机械能，令
\(C=2\mathcal E/r_0^2\)，则 \[
\dot a^2=\frac{K}{a}+C,
\qquad
H^2=Ka^{-3}+Ca^{-2}.
\] 当 \(C>0\)，尺度因子很大时 \(\dot a\) 趋于常数，继续膨胀；当
\(C<0\)，存在 \(\dot a=0\) 的有限最大尺度；当 \(C=0\)，则为本题的
\(t^{2/3}\) 支路。不能单凭 \(\ddot a<0\) 推出必然坍缩。

\textbf{任意初始密度下的年龄积分。} 设当前
\(a=1\)，\(H=H_0>0\)，密度参数为 \(\Omega\)。无压物质满足 \[
\dot a^2=H_0^2\left(\frac{\Omega}{a}+1-\Omega\right).
\] 从膨胀起点积分到当前时刻， \[
H_0t_0=
\int_0^1\sqrt{\frac{a}{\Omega+(1-\Omega)a}}\,da.
\] 这才是非临界情形的年龄；不能把 \(2/(3H_0)\) 用于任意 \(\Omega\)。当
\(\Omega=1\)， \[
H_0t_0=\int_0^1a^{1/2}da=\frac23.
\] 当 \(0<\Omega<1\)，令 \(A=1-\Omega\)，取 \[
a=\frac{\Omega}{A}\sinh^2 u.
\] 则 \[
H_0t=
\frac{\Omega}{A^{3/2}}
\left(\sinh u\cosh u-u\right).
\] 当 \(a=1\)，\(u=\operatorname{arsinh}\sqrt{A/\Omega}\)，得到 \[
\boxed{H_0t_0=
\frac1{1-\Omega}
-\frac{\Omega}{(1-\Omega)^{3/2}}
\operatorname{arsinh}\sqrt{\frac{1-\Omega}{\Omega}}}.
\] 例如 \(\Omega=0.5\)，年龄参数约为 \(0.75355\)，大于临界值
\(2/3\)。在同一个当前 \(H_0\)
下，低密度模型过去的引力减速较弱，回溯到起点所需时间较长。

当 \(\Omega>1\)，令 \(B=\Omega-1\)，用 \[
a=\frac{\Omega}{B}\sin^2 u
\] 得 \[
H_0t=
\frac{\Omega}{B^{3/2}}
\left(u-\sin u\cos u\right).
\] 取 \(a=1\) 后， \[
\boxed{H_0t_0=
\frac{\Omega}{(\Omega-1)^{3/2}}
\arcsin\sqrt{\frac{\Omega-1}{\Omega}}
-\frac1{\Omega-1}}.
\] 例如 \(\Omega=2\)，\(H_0t_0=\pi/2-1\approx0.57080\)。两式在
\(\Omega\to1\) 都趋于
\(2/3\)，但直接数值相减可能丢失精度；接近1时用原积分或级数更稳妥。

\textbf{封闭模型的最大尺度与回落时刻。} 当 \(\Omega>1\)，转折条件
\(\dot a=0\) 给出 \[
a_{\max}=\frac{\Omega}{\Omega-1}.
\] 它对应 \(u=\pi/2\)，从起点到最大尺度的时间为 \[
H_0t_{\max}=
\frac{\pi\Omega}{2(\Omega-1)^{3/2}}.
\] 在当前之后到达最大尺度的等待时间是 \(t_{\max}-t_0\)，而不是
\(t_{\max}\)。理想无压模型的收缩解关于转折时刻对称，总寿命为
\(2t_{\max}\)。这里不包含热压、星系结构或其他成分；它只是共动均匀球的数学延拓。

\needspace{26\baselineskip}

\textbf{当前减速度的无量纲表达。} 定义 \[
q=-\frac{a\ddot a}{\dot a^2}.
\] 在当前，由 \(\ddot a=-4\pi G\rho_0/3\) 与临界密度定义， \[
\boxed{q_0=\Omega/2}.
\] 临界无压解的 \(q=1/2\) 在所有时刻保持不变。把 \(a\propto t^{2/3}\)
代入也有 \[
q=-\frac{(2/3)(-1/3)}{(2/3)^2}=\frac12.
\] 注意 \(\dot a>0\) 与 \(\ddot a<0\)
可同时成立，膨胀与加速膨胀是不同判断。

\textbf{短时间观测量的展开。} 当前后短时间 \(\Delta t\) 内， \[
a(t_0+\Delta t)=
1+H_0\Delta t-\frac{\Omega H_0^2}{4}(\Delta t)^2
+O((H_0\Delta t)^3).
\] 共动间距 \(r=r_0a\) 具有同样的相对展开。一次距离---速度测量确定
\(H_0\)，二阶变化才包含密度参数；只知道线性关系
\(v=H_0r\)，不能唯一确定过去的尺度因子。原题用零能量与质量守恒补足了这一信息。

\subsection{六、关键洞察}\label{ux516dux5173ux952eux6d1eux5bdf-2}

题目的七问是一条连续推导：运动学定义、引力、能量、质量守恒、时间解、单位换算、密度。每一步使用的假设都不同。把
\(H^2=8\pi G\rho/3\) 当作对任意宇宙都成立，会跳过零能量这一必要条件。

共动球的任意半径 \(r_0\)
在最终方程中消去，说明解依赖均匀密度与整体尺度，而不依赖观察者选择多大的球。

\subsection{七、易错点}\label{ux4e03ux6613ux9519ux70b9-2}

把 \(a\) 赋予长度单位；将固定共动标签 \(r_0\)
对时间求导；用球外质量再多加一项引力；把 \(\rho\)
当成恒定而不守恒质量；把 \(H_0^{-1}\)
直接当年龄；把秒差距与百万秒差距少换
\(10^6\)；把题设模型数值写成现实宇宙测量结论。

\subsection{八、课堂讲解建议}\label{ux516bux8bfeux5802ux8bb2ux89e3ux5efaux8bae-2}

每得到一个式子，就在旁边注明它来自哪项假设。图3-2的三条曲线分别画
\(a/a_0\)、\(H/H_0\) 与
\(\rho/\rho_0\)，避免把量纲不同的量放在同一纵轴而不归一化。

让学生检查 \(t=2t_0\)：尺度为 \(2^{2/3}\)，哈勃变量为 \(H_0/2\)，密度为
\(\rho_0/4\)。三个数不应被误写成同一个比例；它们来自不同的时间幂次。

\textbf{用无量纲减速参数复核幂律。} 定义
\(q=-a\ddot a/\dot a^2\)。对临界无压模型 \(a\propto t^{2/3}\)， \[
\dot a=\frac{2a}{3t},\qquad
\ddot a=-\frac{2a}{9t^2},
\qquad \boxed{q=\frac12}.
\] 另外 \[
\dot H=-\frac{2}{3t^2}=-\frac32H^2,
\qquad \frac{\dot\rho}{\rho}=-3H.
\]
把这些关系代入加速度方程和质量守恒式，可同时核对时间幂次与系数。若采用指数膨胀，\(q=-1\)，它不满足本题仅含物质引力的条件。不能因两题都出现
\(H\) 就把第（3）题的幂律与第（5）题规定的指数膨胀混用。

\subsection{九、变式训练}\label{ux4e5dux53d8ux5f0fux8badux7ec3-2}

\textbf{变式一：半密度时刻。} 当密度降为当前的 \(1/2\)
时，求时间、尺度和哈勃变量。

由 \(\rho/\rho_0=(t_0/t)^2\)， \[
\boxed{t=\sqrt2t_0},
\qquad
\boxed{a=2^{1/3}},
\qquad
\boxed{H=H_0/\sqrt2}.
\] 也可先用 \(\rho a^3=\rho_0\) 求尺度，再由 \(a=(t/t_0)^{2/3}\)
求时间。

\textbf{变式二：当前密度偏离临界。} 在当前 \(a=1\)、\(H=H_0\)，密度为
\(\rho_0=\Omega\rho_c\)，其中
\(\rho_c=3H_0^2/(8\pi G)\)。只考虑无压物质，求当 \(\Omega>1\)
时的最大尺度。

能量常量为 \[
C=H_0^2(1-\Omega),
\qquad K=\Omega H_0^2.
\] 转向点 \(\dot a=0\) 满足 \[
\frac{\Omega H_0^2}{a_{\max}}+H_0^2(1-\Omega)=0.
\] 所以 \[
\boxed{a_{\max}=\frac{\Omega}{\Omega-1}}.
\] \(\Omega\downarrow1\)
时最大尺度发散，回到临界模型；这个结论依然限于无压、无宇宙常量的题设框架。

\textbf{从初始数据判断可用的方程。} 引力加速度式与质量守恒共同决定
\(\ddot a\)，还需要初始 \(\dot a\) 才能选择能量常数。可让学生把相同
\(\rho_0\)、不同 \(H_0\)
的三条解分类：机械能为负时回落，为零时临界，为正时持续膨胀。不能仅因``球体均匀''就删掉能量积分常数。

\textbf{在临界点附近检查年龄。} 令 \(\Omega=1+\epsilon\)，从原积分展开：
\[
[\Omega+(1-\Omega)a]^{-1/2}
=[1+\epsilon(1-a)]^{-1/2}.
\] 于是 \[
H_0t_0
=\frac23-\frac{\epsilon}{2}
\int_0^1a^{1/2}(1-a)\,da+O(\epsilon^2)
=\boxed{\frac23-\frac{2\epsilon}{15}+O(\epsilon^2)}.
\]
该结果可以检查两侧闭式公式的极限与符号。密度增加使年龄减小；若展开结果相反，通常是能量常数或
\(1-\Omega\) 写反。

\textbf{换成别的物质成分需要额外物理输入。} 相对论均匀模型中，若压力满足
\(P=w\rho c^2\)，能量守恒给 \[
\rho\propto a^{-3(1+w)}.
\] 在平直、单成分条件下， \[
a\propto t^{2/[3(1+w)]},\qquad
q=(1+3w)/2.
\] 无压物质 \(w=0\) 回到本题，辐射 \(w=1/3\) 给
\(a\propto t^{1/2}\)。这些是增加压力引力与相对论能量守恒后的扩展，不能仅把牛顿质量密度换成辐射密度，就声称原牛顿推导已经包含了辐射宇宙。

\needspace{14\baselineskip}

\textbf{年龄并非仅由 \(H_0\) 决定。} 临界尘埃模型给
\(t_0=2/(3H_0)\)。不同物质密度参数的尘埃模型，年龄积分 \[
H_0t_0=\int_0^1\frac{\sqrt a\,da}{\sqrt{\Omega_{m0}+(1-\Omega_{m0})a}}.
\] 即使当前 \(H_0\)
相同，分母不同也使年龄不同。因此本题的年龄数值由零总能量假设共同决定，不能把
\(2/(3H_0)\) 当成任意宇宙的通式。

密度、标度因子与 \(H\)
的时间幂次要一起检查：\(\rho\propto a^{-3}\)、\(a\propto t^{2/3}\)，推出
\(\rho\propto t^{-2}\)；再用 \(H=2/(3t)\)，推出 Friedmann 比例
\(H^2\propto\rho\)。四项相互一致，比只核验年龄一个数值更可靠。

\newpage

\section{第 4 题}\label{ux7b2c-4-ux9898}

无自旋、无电荷黑洞的史瓦西半径为
\(R_S=2GM/c^2\)。设质子在黑洞外绕转，按类比波尔原子的简化模型，把它称为黑洞---质子原子（BHPA）。

（1）求这种模型中最小可能的轨道半径 \(r_B\)。

（2）求能组成这种模型的黑洞质量上限。

（3）求基态电离能。

（4）题给黑洞辐射温度估算式
\(T_{\rm BH}=10^{23}\,\mathrm{K\cdot kg}/M\)。结合计算数据，讨论上述模型是否可靠。

\includegraphics[width=0.85\linewidth,height=\textheight,keepaspectratio,alt={题4　波尔轨道、视界与模型边界}]{../figures/question-04.pdf}

\subsection{一、题型定位}\label{ux4e00ux9898ux578bux5b9aux4f4d-3}

这道题把牛顿引力、波尔角动量条件和黑洞视界半径放在同一个估算模型中。前三问可以给出形式结果，但靠近视界时牛顿力学和非相对论动能不再可靠；第四问需要据此评估结果，而不能把形式上的极限当成真实稳定原子的证明。

\subsection{二、知识点}\label{ux4e8cux77e5ux8bc6ux70b9-3}

牛顿圆轨道；角动量量子化
\(m_pvr=n\hbar\)；引力束缚能；严格半径条件的上下确界；热辐射能标
\(k_BT\)；无量纲引力耦合；近似条件的相容性。

\subsection{三、物理图像}\label{ux4e09ux7269ux7406ux56feux50cf-3}

图4-1同时画出黑洞视界和模型中的基态半径。增大黑洞质量会使视界半径增大，却使波尔半径减小；两条曲线相交给出形式边界。相交只意味着两个长度相同，并不意味着在那里存在稳定圆轨道。

\subsection{四、解题入口}\label{ux56dbux89e3ux9898ux5165ux53e3-3}

先在远离黑洞的牛顿近似中求 \(r_n(M)\) 与
\(E_n(M)\)，随后检查基态是否位于 \(R_S\)
外。最后比较辐射温度与束缚能，并检验所用速度是否仍远小于光速。

\subsection{五、解析}\label{ux4e94ux89e3ux6790-3}

\textbf{牛顿引力下的波尔轨道。} 记质子质量为
\(m_p\)，忽略质子对黑洞运动的影响，要求 \(M\gg m_p\)。圆轨道满足 \[
\frac{m_pv_n^2}{r_n}=\frac{GMm_p}{r_n^2},
\] 以及类比波尔模型的角动量条件 \[
m_pv_nr_n=n\hbar,\qquad n=1,2,\ldots .
\] 第二式给出 \(v_n=n\hbar/(m_pr_n)\)。代回第一式： \[
\frac{n^2\hbar^2}{m_p^2r_n^2}=\frac{GM}{r_n}.
\] 因此 \[
\boxed{r_n=\frac{n^2\hbar^2}{GMm_p^2}},
\qquad
\boxed{v_n=\frac{GMm_p}{n\hbar}}.
\] 基态半径记为 \[
r_1(M)=\frac{\hbar^2}{GMm_p^2}.
\] 这里采用的是非相对论模型的 \(n=1\)
支路；不能把它直接视为史瓦西时空中的精确量子态。

\textbf{视界外条件与质量上限。} 要求基态在视界之外： \[
r_1>R_S.
\] 即 \[
\frac{\hbar^2}{GMm_p^2}>\frac{2GM}{c^2}.
\] 两边均为正，整理得 \[
M^2<\frac{\hbar^2c^2}{2G^2m_p^2}.
\] 所以形式上 \[
\boxed{M<M_*},
\qquad
\boxed{M_*:=\frac{\hbar c}{\sqrt2Gm_p}}.
\] \(M_*\)
是允许质量集合的上确界。若坚持严格在视界外，等号不是一个允许的稳定状态。

对应的半径下确界为 \[
\begin{aligned}
r_*&=r_1(M_*)
=\frac{\hbar^2}{Gm_p^2}
\frac{\sqrt2Gm_p}{\hbar c}\\
&=\boxed{\frac{\sqrt2\hbar}{m_pc}}.
\end{aligned}
\] 在这一简化模型内，\(r_1>r_*\)，半径可以趋近
\(r_*\)，但不能把视界上的等号写成实际可取的最小轨道。

\textbf{数值量级。} 采用 \[
\hbar=1.05457\times10^{-34}\,\mathrm{J\,s},\qquad
m_p=1.67262\times10^{-27}\,\mathrm{kg},
\] \[
c=2.99792\times10^8\,\mathrm{m/s},\qquad
G=6.67430\times10^{-11}\,\mathrm{m^3\,kg^{-1}\,s^{-2}},
\] 得到 \[
\boxed{M_*\simeq2.00\times10^{11}\,\mathrm{kg}},
\qquad
\boxed{r_*\simeq2.97\times10^{-16}\,\mathrm m}.
\] 这些常量用于补充数值计算，原题的符号表达式不依赖选用的末位数字。

\begin{figure}[H]
\centering
\includegraphics[width=0.817\linewidth,height=\textheight,keepaspectratio,alt={图4-2　半径交点与辐射---束缚能标比较}]{../figures/solution-04.pdf}
\caption*{图4-2　半径交点与辐射---束缚能标比较}
\end{figure}

\textbf{基态电离能。} 牛顿圆轨道中 \[
K_n=\frac12m_pv_n^2=\frac{GMm_p}{2r_n},
\] 而引力势能 \[
U_n=-\frac{GMm_p}{r_n}.
\] 因此总能量 \[
\boxed{E_n=-\frac{G^2M^2m_p^3}{2n^2\hbar^2}}.
\] 从基态提升到无穷远、末速度为零，所需能量为 \[
\boxed{E_{\rm ion}(M)=\frac{G^2M^2m_p^3}{2\hbar^2}}.
\] 质量 \(M\) 越大，模型束缚越强。若将其形式外推到 \(M_*\)， \[
E_{\rm ion,*}=\frac14m_pc^2.
\] 所以其上确界为 \[
\boxed{E_{\rm ion,*}\simeq3.76\times10^{-11}\,\mathrm J
\simeq235\,\mathrm{MeV}}.
\] ``基态电离能''一般是 \(M\) 的函数。\(m_pc^2/4\)
仅是沿前三问形式边界得到的极限值，不能对任意黑洞质量直接使用。

\textbf{辐射温度的量级。} 按原题的近似温度公式，在质量上界附近 \[
T_*\simeq\frac{10^{23}}{2.00253\times10^{11}}\,\mathrm K
\simeq4.99\times10^{11}\,\mathrm K.
\] 以 \(k_B=1.380649\times10^{-23}\,\mathrm{J/K}\)，热能标为 \[
\boxed{k_BT_*\simeq6.89\times10^{-12}\,\mathrm J
\simeq43.0\,\mathrm{MeV}}.
\] 故 \[
\boxed{\frac{k_BT_*}{E_{\rm ion,*}}\simeq0.183}.
\] 它小于1，但不小到可忽略；黑体光谱也不是以 \(k_BT\)
为绝对截止。仅凭这一比值，不能断言没有高能辐射、不会发生电离，或能形成长寿命束缚态。寿命问题还需要辐射谱、吸收截面、黑洞质量变化及量子态的吸收边界条件。

\textbf{减小质量会怎样。} 令 \(x=M/M_*\)。形式模型给出 \[
r_1=\frac{r_*}{x},\qquad
R_S=xr_*,\qquad
E_{\rm ion}=x^2E_{\rm ion,*},\qquad
T_{\rm BH}=\frac{T_*}{x}.
\] 因而 \[
\boxed{\frac{k_BT_{\rm BH}}{E_{\rm ion}}
\simeq\frac{0.183}{x^3}}.
\] 当 \(x<0.183^{1/3}\simeq0.568\)，热能标就超过形式束缚能。例如
\(x=0.1\)
时，轨道距视界较远，但比值已约为183；减小质量虽然改善弱场近似，却使热辐射相对束缚更加显著。

\textbf{非相对论近似在上界处已经失效。} 引入 \[
\alpha_G=\frac{GMm_p}{\hbar c}.
\] 模型中的基态速度为 \(v_1=\alpha_Gc\)，半径比为 \[
\frac{r_1}{R_S}=\frac{1}{2\alpha_G^2}.
\] 在 \(M=M_*\)， \[
\alpha_G=\frac1{\sqrt2},\qquad v_1=\frac c{\sqrt2}.
\] 如此大的速度不能用 \(K=m_pv^2/2\)
精确描述，且轨道已处在强引力区域。因此前三问算出的极限本身正好位于模型失效处。可靠的牛顿、非相对论处理需要
\(\alpha_G\ll1\)，比``仅在视界外''严格得多。

\textbf{用霍金精确系数检查参数相容性。} 无自旋、无电荷黑洞的温度是 \[
T_H=\frac{\hbar c^3}{8\pi G M k_B}.
\] 用它与同一牛顿束缚能比较： \[
\frac{k_BT_H}{E_{\rm ion}}
=\frac{1}{4\pi\alpha_G^3}.
\] 当 \(\alpha_G\ll1\)
时，这个比值远大于1。由此可见，既要求牛顿近似很准、又要求热辐射能标远低于束缚能，无法在这套简化模型中同时实现。此比较是在评估近似的相容性，并不是完成了真实黑洞量子束缚态的寿命计算。温度公式可参看
\href{https://astrofrelat.fcaglp.unlp.edu.ar/agujeros_negros/media/Papers/Hawking_1975-Particle_creation_by_black_holes.pdf}{Hawking
原始论文}。

\textbf{视界外也不等于稳定圆轨道。} 史瓦西时空中的经典稳定类时圆轨道要求
\(r>6GM/c^2=3R_S\)；边界是边际稳定轨道。仅用 \(r>R_S\)
检查会把视界外的一部分不稳定或不存在的圆轨道也算进去。关于圆轨道有效势的推导，可参看
\href{https://arxiv.org/pdf/1406.4282}{d'Ambrosi 与 van Holten
的原始论文，第2节}。

这条经典结论不能简单替换成完整量子基态条件，但足以指出题目混合模型的缺项。最终评估是：前三问给出牛顿---波尔估算和形式上下界；给定辐射数据与近似检查不足以支持真实、稳定、长寿命
BHPA 的存在。

\textbf{把``轨道在外面''拆成三个不同条件。} 用长度 \(r_g=GM/c^2\)
表示黑洞尺度，视界是
\(2r_g\)。经典史瓦西圆轨道中，以单位质子质量的角动量 \(\ell\)
为变量，径向有效势可写为 \[
V_{\rm eff}(r)=
\left(1-\frac{2r_g}{r}\right)
\left(c^2+\frac{\ell^2}{r^2}\right).
\] 圆轨道要求 \(V_{\rm eff}'=0\)： \[
\frac{2r_gc^2}{r^2}
-\frac{2\ell^2}{r^3}
+\frac{6r_g\ell^2}{r^4}=0.
\] 解出 \[
\boxed{\ell^2=\frac{r_gc^2r^2}{r-3r_g}}.
\] 角动量平方为正先要求 \(r>3r_g\)，比视界外条件 \(r>2r_g\) 更强。再固定
\(\ell\) 对有效势求二阶导，并代入圆轨道关系： \[
V_{\rm eff}''=
\frac{2r_gc^2(r-6r_g)}
{r^3(r-3r_g)}.
\] 所以稳定要求 \(r>6r_g=3R_S\)；\(3r_g<r<6r_g\)
的类时圆轨道虽存在，却不稳定；\(2r_g<r\le3r_g\)
不能成为这种质点的圆轨道。这个有效势检查说明，原题的视界条件只能给形式估算上界。

\textbf{在圆轨道关系中插入角动量量子化会怎样。}
此处仅作半经典估算，不代替黑洞外完整波动方程。若仍令
\(m_p\ell=n\hbar\)，记 \[
x=r/r_g,\qquad j=n/\alpha_G,
\] 则 \[
j^2=\frac{x^2}{x-3}.
\] 它化为二次方程 \[
x^2-j^2x+3j^2=0.
\] 两根为 \[
x_\pm=\frac{j^2\pm j\sqrt{j^2-12}}2.
\] 实圆轨道先要求 \(j\ge\sqrt{12}\)，即 \[
\alpha_G\le n/\sqrt{12}.
\] 其中 \(x_+\ge6\) 是稳定支，\(x_-\le6\) 是不稳定支；等号在 \(x=6\)
合并。牛顿波尔模型把两支的结构全部省略了。在弱场 \(j\gg1\)， \[
x_+=j^2-3-\frac9{j^2}+O(j^{-4}),
\] 对应牛顿半径 \(r_n/r_g=n^2/\alpha_G^2\) 的首项。弱场展开的适用条件是
\(\alpha_G/n\ll1\)，不能延伸到两根合并处。

\textbf{局部速度与能量也显示模型边界。}
史瓦西圆轨道对当地静止观察者的速度为 \[
\frac{v_{\rm local}^2}{c^2}
=\frac{r_g}{r-2r_g}.
\] 在最内稳定圆轨道 \(r=6r_g\)， \[
v_{\rm local}=c/2.
\] 每单位静质量的轨道总能量为 \[
\frac{\mathcal E}{m_pc^2}
=\frac{1-2r_g/r}{\sqrt{1-3r_g/r}}.
\] 在 \(r=6r_g\)，其值为 \(\sqrt{8/9}\)，相对无穷远静止质子的束缚比例为
\[
1-\sqrt{8/9}\approx0.05719.
\] 它与形式牛顿视界上界处的 \(1/4\)
不同，因为两个计算采用的半径条件和动力学均不同。不能把前者当成原题牛顿模型的``精确数值修正''直接替换；应说明模型已改变。

\needspace{14\baselineskip}

\textbf{数值比较应在共同有效的范围内进行。} 取 \(\alpha_G=0.05,n=1\)，则
\(j=20\)。稳定圆轨道的无量纲半径为 \[
x_+=\frac{400+\sqrt{155200}}2\approx396.98.
\] 牛顿半径为 \(400r_g\)，相对差约
\(0.76\%\)。此时两种几何估算接近，但霍金热能仍远大于牛顿束缚能：热效应与轨道几何误差是两个独立检验，不能因半径一致就断言束缚态长期可靠。

\textbf{有限黑洞质量带来的约化质量修正。}
若暂不讨论视界，只修正牛顿两体运动，约化质量为 \[
\mu=\frac{Mm_p}{M+m_p}.
\] 相对坐标的波尔半径改为 \[
r_n=\frac{n^2\hbar^2}{\mu GMm_p},
\] 能级为 \[
E_n=-\frac{\mu G^2M^2m_p^2}{2n^2\hbar^2}.
\]
在本题黑洞质量数量级下，\(\mu/m_p=1-O(m_p/M)\)，修正极小。保留约化质量却忽略相对论与辐射，并不能提高主要结论的可靠性；应先比较各项误差的尺度。

\subsection{六、关键洞察}\label{ux516dux5173ux952eux6d1eux5bdf-3}

一个形式上能算出的极值，可能恰好位于近似失效的区域。这里半径最小、质量最大和束缚最强同时把轨道推近视界，也同时把速度推到相对论量级。第四问应检查这条逻辑链，而不是只比较两个大数。

热能标的意义是光谱量级，不是每个辐射粒子的固定能量。\(k_BT<E_{\rm ion}\)
不能直接等价为``不会电离''。

\subsection{七、易错点}\label{ux4e03ux6613ux9519ux70b9-3}

把 \(h\) 代替 \(\hbar\) 而多出
\(2\pi\)；忘记电离能是负总能量的绝对值；将严格不等式的上确界写成允许质量；把视界半径当作最内稳定轨道；在
\(v\sim c\) 时仍宣称牛顿值精确；用单个温度比值代替寿命计算。

\subsection{八、课堂讲解建议}\label{ux516bux8bfeux5802ux8bb2ux89e3ux5efaux8bae-3}

让学生先在图4-2上比较 \(r_1\propto M^{-1}\) 与
\(R_S\propto M\)，再算交点。完成前三问后，立即计算对应速度，使模型局限成为推导的一部分。

可以用 \(M/M_*=1,0.5,0.1\)
做表格：半径、视界和束缚能按不同幂次改变。分别检查``外部轨道''\,``弱场''\,``非相对论''和``辐射扰动小''，避免把四个条件压缩成一个
\(r>R_S\)。

\textbf{Schwarzschild 圆轨道的能量。} 设
\(r_g=GM/c^2\)，对有质量粒子的圆轨道，以静质量能量为单位的守恒能量满足
\[
\mathcal E^2=\frac{(1-2r_g/r)^2}{1-3r_g/r}.
\] 在最内稳定圆轨道 \(r=6r_g\)， \[
\mathcal E=\sqrt{8/9},\qquad
1-\mathcal E\simeq0.05719.
\]
这可由有效势驻值条件消去角动量得到。它是经典相对论轨道的结合能比例，不是本题半经典轨道量子化已经给出的普遍辐射效率。把这项检查放在教师补充中，可进一步解释为什么牛顿公式不能被直接外推到黑洞附近。

\subsection{九、变式训练}\label{ux4e5dux53d8ux5f0fux8badux7ec3-3}

\textbf{变式一：第 \(n\) 层的视界外条件。} 由 \(r_n=n^2r_1\)，要求
\(r_n>R_S\) 得 \[
\boxed{M<nM_*}.
\] 对应形式半径下确界为 \(nr_*\)，电离能极限仍为
\(m_pc^2/4\)。这是因为交点处 \(v_n=c/\sqrt2\)，模型在每个 \(n\)
的边界都有同样的失效问题。

\textbf{变式二：在弱场参数下比较能标。} 取
\(\alpha_G=0.05\)，不计算任何量纲常量，求半径比、速度比及霍金热能标与基态束缚能之比。
\[
\boxed{r_1/R_S=200},\qquad
\boxed{v_1/c=0.05},
\] \[
\boxed{k_BT_H/E_{\rm ion}=\frac{1}{4\pi(0.05)^3}\simeq637}.
\]
轨道已处在较好的弱场与低速区域，但热辐射相对束缚仍很强。该例直接说明改善一个近似条件不会自动改善所有物理条件。

\textbf{半经典稳定轨道估算的边界。}
上述角动量插入法说明牛顿形式上界为何过宽，但它没有证明存在一个精确的``质子基态''。完整问题还涉及视界的吸收边界、波函数衰减以及黑洞辐射背景。评分时可认可学生用
\(\alpha_G\) 和 \(r/R_S\)
证明弱场失效，不要求把额外半经典圆轨道当作原题唯一解。

\textbf{温度公式的系数与单位。} 题给 \(10^{23}\,\mathrm{K\,kg}/M\)
是数量级式。精确霍金温度的系数为 \[
\frac{\hbar c^3}{8\pi Gk_B}.
\] 讨论电离时比较的是 \(k_BT\) 与
\(|E_1|\)，不能直接比较开尔文和焦耳。将精确式除以牛顿束缚能，量纲常量消去：
\[
\frac{k_BT_H}{|E_1|}
=\frac1{4\pi\alpha_G^3}.
\] 若要求该比例远小于1，则 \(\alpha_G\) 需要大；若要求非相对论束缚，则
\(\alpha_G\ll1\)。两者没有共同的参数极限。这个矛盾比单个温度的数值更能说明模型不足。

\textbf{对高激发态的扩展。} 第 \(n\) 级有 \(|E_n|=|E_1|/n^2\)，速度尺度
\(\alpha_Gc/n\)。固定黑洞质量时， \[
\frac{k_BT_H}{|E_n|}
=\frac{n^2}{4\pi\alpha_G^3}.
\] 提高 \(n\)
可以减小轨道速度，却进一步削弱束缚相对热能。不能用高激发态同时解决非相对论和热稳定两个问题。若学生提出``把质子放到更远的轨道就安全''，可以用这两个比例要求其逐项检查。

\needspace{14\baselineskip}

\textbf{阈值附近的近似层级。} Newton--Bohr 轨道半径与 Schwarzschild
半径相等时，轨道速率已达
\(c/\sqrt2\)。这正是非相对论动能展开不受控的区域，所以该质量只是形式边界。不能一边用非相对论轨道，一边把等号点称为可靠的最小稳定原子。

一般相对论的圆轨道稳定半径为
\(r>6GM/c^2\)，不同于只要求轨道在视界外。把角动量半经典地设为 \(n\hbar\)
后得到的圆轨道根，也不是完整量子束缚谱。正文给这个补充只是展示稳定条件改变，而不是用一个额外方程掩盖旧模型的适用范围。

温度比较还须固定采用原题 \(T=10^{23}/M\) 的近似系数，还是 Hawking
的精确物理常数式。不同口径的数值不能在同一不等式里混搭。本批保留原题数值模型，并独立说明量级与物理局限。

\newpage

\section{第 5 题}\label{ux7b2c-5-ux9898}

一团平衡辐射气体体积为 \(V\)，单位体积内能为 \(u=\beta T^4\)，压强为
\(P=u/3\)，其中 \(\beta\) 为常量，\(T\) 为热力学温度。

（1）用热力学第一定律求 \(\delta Q\)，以 \(V,T\) 为变量。

（2）由 \(dS=\delta Q_{\rm rev}/T\) 求熵。

（3）把它放在绝热膨胀的球形容器中，半径满足
\(\dot r=\alpha r\)，\(\alpha\) 为常量，初态 \(r=r_0,T=T_0\)。求
\(T(t)\)。

（4）给定
\(\alpha=72\,\mathrm{km\,s^{-1}\,Mpc^{-1}}\)，\(1\,\mathrm{Mpc}=3.26\times10^6\,\mathrm{ly}\)，求
\(\alpha\) 的量纲与SI数值。

\includegraphics[width=0.84\linewidth,height=\textheight,keepaspectratio,alt={题5　辐射气体绝热膨胀}]{../figures/question-05.pdf}

\subsection{一、题型定位}\label{ux4e00ux9898ux578bux5b9aux4f4d-4}

以 \(V,T\)
为状态变量求辐射气体的热量微分与熵，再把绝热关系和规定的半径增长率组合起来。体积内能
\(u\) 与总内能 \(U\) 的区别决定了求微分时必须保留 \(dV\) 项。

\subsection{二、知识点}\label{ux4e8cux77e5ux8bc6ux70b9-4}

总内能 \(U=uV\)；第一定律
\(\delta Q=dU+P\,dV\)；可逆熵微分；微分的积分性；绝热不变量；指数增长；逆时间量纲。

\subsection{三、物理图像}\label{ux4e09ux7269ux7406ux56feux50cf-4}

辐射气体随膨胀对容器做功，温度下降，单位体积内能按 \(T^4\)
下降。总内能还乘以增加的体积，故它的变化不能只从 \(u\) 判断。图5-2把
\(r,V,T,U\) 分别按各自初值归一化，显示它们的不同指数。

\subsection{四、解题入口}\label{ux56dbux89e3ux9898ux5165ux53e3-4}

先写 \(U=\beta VT^4\) 并全微分，再加压强功。除以 \(T\)
后检查两个系数能否组成同一个状态函数；求熵后，绝热关系可直接由 \(S\)
不变得到。

\subsection{五、解析}\label{ux4e94ux89e3ux6790-4}

\textbf{第一定律的完整微分。} 总内能 \[
U=\beta VT^4.
\] 在 \(V,T\) 都可能变化时， \[
dU=\beta T^4\,dV+4\beta VT^3\,dT.
\] 压力为 \[
P=\frac13\beta T^4.
\] 以系统吸热为正、系统对外做功为 \(P\,dV\)，第一定律为 \[
\delta Q=dU+P\,dV.
\] 代入得 \[
\boxed{\delta Q=4\beta VT^3\,dT+\frac43\beta T^4\,dV}.
\]
第一项是温度改变的内能贡献；第二项包含新增体积的内能与膨胀功。不能只写
\(\delta Q=V\,du+P\,dV\) 而漏掉 \(u\,dV\)。

\textbf{熵微分。} 对可逆准静态过程， \[
dS=\frac{\delta Q_{\rm rev}}T
=4\beta VT^2\,dT+\frac43\beta T^3\,dV.
\] 右侧恰好是 \[
d\left(\frac43\beta VT^3\right).
\] 所以 \[
\boxed{S(V,T)=\frac43\beta VT^3+S_*}.
\] 常量 \(S_*\) 不影响熵差。如果以无辐射的零温态为熵零点，可以取
\(S_*=0\)。这里只依据题设辐射模型确定状态依赖，不能从微分式单独确定任意系统的绝对熵常量。

\textbf{积分性的独立检查。} 若写 \[
dS=A(V,T)\,dT+B(V,T)\,dV,
\] 则存在状态函数的必要条件为 \[
\left(\frac{\partial A}{\partial V}\right)_T
=\left(\frac{\partial B}{\partial T}\right)_V.
\] 本题 \[
A=4\beta VT^2,\qquad B=\frac43\beta T^3,
\] 因此两侧均为
\(4\beta T^2\)。这一步检查了压力与内能公式在热力学上相容，而不仅是``能把式子积分出来''。

\begin{figure}[H]
\centering
\includegraphics[width=0.648\linewidth,height=\textheight,keepaspectratio,alt={图5-2　绝热不变量与归一化时间曲线}]{../figures/solution-05.pdf}
\caption*{图5-2　绝热不变量与归一化时间曲线}
\end{figure}

\textbf{绝热关系。} 取无热量交换的准静态膨胀，令 \(\delta Q=0\)： \[
4\beta VT^3\,dT+\frac43\beta T^4\,dV=0.
\] 在 \(V,T>0\) 时除以 \(4\beta VT^4\)： \[
\frac{dT}T+\frac13\frac{dV}V=0.
\] 积分为 \[
\boxed{TV^{1/3}=\text{常量}},
\qquad\boxed{VT^3=\text{常量}}.
\] 第二个形式正对应熵不变。再结合 \(P\propto T^4\)： \[
\boxed{PV^{4/3}=\text{常量}}.
\] 该指数来自辐射状态关系，不应借用单原子理想气体的 \(5/3\)。

\textbf{球形容器的温度。} 对球体 \[
V=\frac{4\pi}{3}r^3,
\qquad V^{1/3}\propto r.
\] 因此绝热不变量等价于 \[
\boxed{Tr=T_0r_0}.
\] 规定的半径方程为 \(\dot r=\alpha r\)，分离变量得 \[
\frac{dr}r=\alpha\,dt,
\] 结合 \(r(0)=r_0\)： \[
\boxed{r(t)=r_0e^{\alpha t}}.
\] 所以 \[
\boxed{T(t)=T_0e^{-\alpha t}}.
\] 其余量分别为 \[
\frac V{V_0}=e^{3\alpha t},\qquad
\frac u{u_0}=\frac P{P_0}=e^{-4\alpha t},\qquad
\frac U{U_0}=e^{-\alpha t}.
\] 总内能只按 \(e^{-\alpha t}\) 下降，而能量密度按 \(e^{-4\alpha t}\)
下降，差别来自体积的 \(e^{3\alpha t}\) 因子。

\textbf{功率复核。} 对上述解， \[
\dot U=-\alpha U.
\] 另一方面， \[
\dot V=3\alpha V,\qquad P=\frac{U}{3V},
\] 故 \[
P\dot V=\alpha U.
\] 于是 \[
\dot U+P\dot V=0,
\] 正好满足绝热第一定律。若温度指数误写成
\(-3\alpha t\)，这一能量核验会立即失败。

\textbf{\(\alpha\) 的量纲和数值。} 由 \(\dot r=\alpha r\)， \[
[\alpha]=\frac{[\dot r]}{[r]}=\mathrm{s^{-1}}.
\] ``每百万秒差距72千米每秒''是速度除以距离，不是一个单独的速度。采用
\(c=2.99792458\times10^8\,\mathrm{m/s}\)、儒略年 \(365.25\) 天， \[
1\,\mathrm{ly}=c(365.25\times86400\,\mathrm s)
\simeq9.46073\times10^{15}\,\mathrm m.
\] 题给换算得到 \[
1\,\mathrm{Mpc}\simeq
3.26\times10^6\times9.46073\times10^{15}\,\mathrm m
\simeq3.08420\times10^{22}\,\mathrm m.
\] 所以 \[
\boxed{\alpha\simeq2.3345\times10^{-18}\,\mathrm{s^{-1}}}.
\] 若只保留题给 \(3.26\) 的精度，可写
\(2.33\times10^{-18}\,\mathrm{s^{-1}}\)。它与第3题用
\(1\,\mathrm{pc}=3.09\times10^{16}\,\mathrm m\)
换算出的末位略有差别，源于输入近似，不能把两个换算值混用后声称更多有效数字。

\textbf{时间尺度。} 半径增长一倍、温度降为一半所需时间相同： \[
e^{\alpha t_{1/2}}=2,
\qquad
\boxed{t_{1/2}=\frac{\ln2}{\alpha}}.
\] 数值约为 \(2.97\times10^{17}\,\mathrm s\)，即
\(9.41\times10^9\,\mathrm{yr}\)。在一个 \(\alpha^{-1}\) 时间内，半径增为
\(e\) 倍，温度降为 \(e^{-1}\) 倍；这是指数模型的自然时间尺度。

\textbf{与第3题模型的区别。} 本题直接规定 \(\alpha\)
为常量，因此半径指数增长。第3题的物质零能量模型则有
\(H=2/(3t)\)，并非常量。两题使用相同形式的``速度正比距离''，不代表具有相同的动力学或年龄关系。若本题再要求自洽的宇宙动力学，还需给出驱动指数膨胀的条件；现有计算只是题设容器运动下的辐射热学。

\textbf{熵密度和总熵。} 熵密度 \[
s=\frac SV=\frac43\beta T^3
\] 在零常量约定下随时间按 \(e^{-3\alpha t}\) 下降，而总熵 \(S=sV\)
保持常量。密度减小并不表示绝热可逆过程损失了熵，它是同一总熵分布在更大体积中的结果。

\textbf{有限绝热膨胀的功与能量账。} 令体积比 \(z=V_2/V_1>1\)。绝热式
\(VT^3=\mathrm{const}\) 给 \[
T_2=T_1z^{-1/3},\qquad
U_2=U_1z^{-1/3}.
\] 辐射对外做功 \[
W=\int_{V_1}^{V_2}P\,dV
=\frac{\beta T_1^4V_1^{4/3}}3
\int_{V_1}^{V_2}V^{-4/3}dV.
\] 因此 \[
\boxed{W=U_1(1-z^{-1/3})=U_1-U_2}.
\] 与 \(\delta Q=0\)
的第一定律一致。膨胀十倍体积并不使总辐射能降为十分之一；总能比例为
\(10^{-1/3}\)，能量密度比例为
\(10^{-4/3}\)。总能与密度的幂次差来自体积的变化。

\textbf{等温膨胀为什么需要额外补充辐射能。} 等温时 \(U=\beta VT^4\)
随体积增大，内能并不恒定。对 \(\Delta V>0\)， \[
\Delta U=\beta T^4\Delta V,\qquad
W=\frac13\beta T^4\Delta V,
\] \[
Q=\frac43\beta T^4\Delta V.
\]
这是辐射与定量理想气体的差别：定量理想气体内能只随温度，辐射的平衡总能还随容器体积。若把本题套进
\(Q=W\) 的等温气体公式，会漏掉 \(3Q/4\) 对应的内能增长。

\textbf{用辐射为工质的卡诺循环。} 取高、低温 \(T_h>T_c\)。高温等温膨胀由
\((V_a,T_h)\) 到 \((V_b,T_h)\)，然后绝热膨胀到
\((V_c,T_c)\)。低温等温压缩到
\((V_d,T_c)\)，最后绝热压缩返回。两绝热段要求 \[
V_cT_c^3=V_bT_h^3,\qquad
V_dT_c^3=V_aT_h^3.
\] 所以 \[
V_c-V_d=(V_b-V_a)(T_h/T_c)^3.
\] 吸热、放热的大小分别为 \[
Q_h=\frac43\beta T_h^4(V_b-V_a),
\] \[
|Q_c|=\frac43\beta T_c^4(V_c-V_d).
\] 相除后， \[
\boxed{\frac{|Q_c|}{Q_h}=\frac{T_c}{T_h}},
\qquad
\boxed{\eta=1-\frac{T_c}{T_h}}.
\]
这个结果不是事先套用效率公式，而是从辐射状态方程和两段绝热关系算出。它同时验证了熵函数：高温吸热引起
\(\Delta S=Q_h/T_h\)，低温放热恰好减少同样的熵。

\needspace{14\baselineskip}

\textbf{为何不存在任意独立的``辐射定压热容''。} 对平衡辐射
\(P=\beta T^4/3\)，定压即定温。沿等压体积变化可以吸热，但温度不变，\(\delta Q/dT\)
不是通常有限的定压热容。定容热容则明确： \[
C_V=\left(\frac{\partial U}{\partial T}\right)_V=4\beta VT^3.
\] 常见气体公式 \(C_P-C_V=R\)
依赖定量粒子状态方程，不能在没有固定粒子数的辐射中直接使用。类似地，绝热指数
\(4/3\) 可由 \(PV^{4/3}=\mathrm{const}\) 定义，但不应在此不加说明地写成
\(C_P/C_V\)。

\textbf{一般幂律介质的积分性检查。} 为比较模型，设平衡能量密度
\(u=A T^s\)、压力 \(P=w u\)，其中 \(A,w,s\) 为常量。由第一定律， \[
\frac{\delta Q}{T}
=sAVT^{s-2}dT+(1+w)AT^{s-1}dV.
\] 要使它成为熵的恰当微分，需有 \[
sAT^{s-2}=(1+w)(s-1)AT^{s-2}.
\] 因此 \[
\boxed{s=\frac{1+w}{w}}.
\] 取 \(w=1/3\) 恰得 \(s=4\)。若 \(u(T)\) 与 \(P(u)\)
任意选取，第一定律虽然能给一条热量表达式，却未必有状态熵。本题两关系相容，这是推出可逆绝热式与卡诺循环的前提。

\subsection{六、关键洞察}\label{ux516dux5173ux952eux6d1eux5bdf-4}

辐射气体的压力与能量密度之比决定绝热指数。由
\(P=u/3\)、\(u\propto T^4\)，第一定律直接给出 \(Tr\)
不变。指数膨胀只负责把这个不变量转换为时间函数。

除以温度后检查积分性，可帮助辨认一个热学模型是否自洽。它与第045套指定解离度的模型形成有用比较：能量和压力分别有公式，并不自动保证它们共同允许熵状态函数。

\subsection{七、易错点}\label{ux4e03ux6613ux9519ux70b9-4}

把 \(u\) 与 \(U\) 混用；全微分中漏掉
\(u\,dV\)；把不可逆过程的实际热量直接除以温度当作熵差；球体体积指数少乘3；把总熵与熵密度混淆；把常量
\(\alpha\) 的指数模型套入物质宇宙的 \(2/(3H_0)\) 年龄式。

\subsection{八、课堂讲解建议}\label{ux516bux8bfeux5802ux8bb2ux89e3ux5efaux8bae-4}

先写三行：\(U=uV\)、\(dU=u\,dV+V\,du\)、\(\delta Q=dU+P\,dV\)。逐项确认之后再代入
\(u=\beta T^4\)。最后用 \(P\dot V=-\dot U\) 做能量检查。

图5-2应显示四条归一化曲线，而不是把半径、温度和能量放在不注明单位的共同纵轴。让学生预测温度降为一半时：半径为2倍、体积为8倍、能量密度为
\(1/16\)、总内能为 \(1/2\)，再与图对照。

\textbf{可逆绝热与绝热自由膨胀。} 若辐射气体从 \(V_i\) 扩到
\(V_f\)，可逆绝热给 \(TV^{1/3}\)
恒定。另一变式是隔热刚性总容器内移走隔板，不对外做功，最终仍达到题设的平衡辐射态。此时
\(U\) 恒定： \[
T_f=T_i(V_i/V_f)^{1/4}.
\] 最终熵比为 \[
\frac{S_f}{S_i}=(V_f/V_i)^{1/4}>1.
\] 同为``绝热''，自由膨胀不是等熵过程。若 \(V_f=8V_i\)，两种温度比分别为
\(1/2\) 与
\(8^{-1/4}\)。应明确最终平衡假设；不能把任意未经热化的光子分布都当作单温黑体辐射。

\subsection{九、变式训练}\label{ux4e5dux53d8ux5f0fux8badux7ec3-4}

\textbf{变式一：任意给定膨胀率。} 若
\(\dot r=\alpha(t)r\)，仍保持绝热准静态，求温度。 \[
\boxed{T(t)=T_0\exp\left[-\int_0^t\alpha(s)\,ds\right]}.
\] 这个结论只使用 \(Tr\) 不变。当 \(\alpha(t)=b/(t+t_*)\) 时， \[
\boxed{T(t)=T_0\left(\frac{t_*}{t+t_*}\right)^b}.
\]

\textbf{变式二：固定温度膨胀。} 在温度 \(T\) 不变时，体积从 \(V_1\) 增至
\(V_2\)，求吸热、做功和内能变化。 \[
\Delta U=\beta T^4(V_2-V_1),
\] \[
W=\frac13\beta T^4(V_2-V_1),
\] \[
\boxed{Q=\frac43\beta T^4(V_2-V_1)}.
\] 同时 \[
\Delta S=\frac43\beta T^3(V_2-V_1)=Q/T.
\]
辐射气体等温膨胀时总内能增加，不能照搬固定粒子数理想气体的``等温内能不变''。

\textbf{把模型自洽检查加入评讲。}
可与第045套的反应气体题对照：两题都能写 \(dU+P\,dV\)，但只有先检查
\(\delta Q/T\) 的积分性，才知道是否存在一致的熵。这里 \(P=u/3\) 与
\(u\propto T^4\) 恰好满足积分性，因而等温---绝热循环给出卡诺效率。

\textbf{带热源的指数膨胀。} 若仍规定
\(\dot V=3\alpha V\)，并加入恒定功率 \(\mathcal P\)，第一定律成为 \[
\dot U+\alpha U=\mathcal P.
\] 解为 \[
U(t)=\frac{\mathcal P}{\alpha}
+\left(U_0-\frac{\mathcal P}{\alpha}\right)e^{-\alpha t}.
\] 由于 \(V=V_0e^{3\alpha t}\)， \[
T(t)=\left[\frac{U(t)}{\beta V_0e^{3\alpha t}}\right]^{1/4}.
\] 若 \(\mathcal P=0\)，它回到 \(T=T_0e^{-\alpha t}\)；若
\(\mathcal P>0\) 且时间很长，总能趋于 \(\mathcal P/\alpha\)，而温度以
\(e^{-3\alpha t/4}\)
衰减。这说明``温度下降''并不等同``总能下降''，还要查体积与热输入。

原题是准静态热学模型。若膨胀速度已经达到不能维持各处热平衡的范围，则单一温度
\(T(t)\)
的描述需另作检验；这一限制来自热平衡假设，不能通过单纯更换指数系数解决。

\needspace{14\baselineskip}

\textbf{等温辐射循环的热量。} 辐射 \(U=\beta VT^4\)，等温体积增大时
\(dU=\beta T^4dV\)，功为 \(P\,dV=\beta T^4dV/3\)。因此吸热 \[
\delta Q=\frac43\beta T^4dV.
\] 它不是气体等温过程中 \(dU=0\) 的情形。两个辐射绝热段满足 \(TV^{1/3}\)
恒定，由此等温段体积增量的比为 \((T_h/T_c)^3\)，热量比最终仍为
\(T_h/T_c\)，得到 Carnot 效率。

教师可先让学生检查为何辐射等温膨胀仍需要补充能量，再求循环；这能避免把理想气体的内能只依赖温度经验直接迁移到无固定粒子数的光子气体。

\newpage

\section{第 6 题}\label{ux7b2c-6-ux9898}

两面半无限平面镜夹角为 \(72^\circ\)，关于零度方向对称，镜面方向分别为
\(+36^\circ\) 与 \(-36^\circ\)。夹角内放置圆弧形发光标尺，刻度为
\(-36^\circ\) 到 \(+36^\circ\)。观察者在角方向 \(\phi\)，其中
\(-36^\circ<\phi<36^\circ\)。在绕镜面交点的角位置 \(\theta\)
看见的像所对应的原标尺刻度记为 \(\varphi(\theta,\phi)\)。

求以下四个角区间内的 \(\varphi(\theta,\phi)\)，必要时分段：

（1）\(36^\circ<\theta<108^\circ\)；（2）\(108^\circ<\theta<180^\circ\)；（3）\(180^\circ<\theta<252^\circ\)；（4）\(252^\circ<\theta<324^\circ\)。

\includegraphics[width=0.86\linewidth,height=\textheight,keepaspectratio,alt={题6　72度双镜、标尺与角度基准}]{../figures/question-06.pdf}

\subsection{一、题型定位}\label{ux4e00ux9898ux578bux5b9aux4f4d-5}

双镜的反射像可通过把镜面和物体展开到周围扇区求出。半无限镜还要求实际光路在镜面正向延伸部分发生反射。观察者位置影响第二、第三扇区的可见路径，因此不能只写一个与
\(\phi\) 无关的周期折叠式。

\subsection{二、知识点}\label{ux4e8cux77e5ux8bc6ux70b9-5}

关于方向角 \(\gamma\) 的反射
\(\varphi\mapsto2\gamma-\varphi\)；交替镜像；展开与折回；直线连线绕原点的角跨度；镜面交点的排除；分段函数。

\subsection{三、物理图像}\label{ux4e09ux7269ux7406ux56feux50cf-5}

图6-1以镜面交点 \(O\) 为圆心，零度方向为两镜角平分线。\(\theta\)
是像点相对这一基准的方向角，\(\phi\)
是观察者方向。它们都不是观察者眼睛处的视线偏角。

展开光路后，曲折反射路变为观察者到虚像的一条直线。若把同一个周向位置分别记成
\(\theta\) 和
\(\theta-360^\circ\)，便得到沿上镜或下镜展开的两条候选路径；真正可用的是两端绕原点角跨度小于
\(180^\circ\) 的那条。

\subsection{四、解题入口}\label{ux56dbux89e3ux9898ux5165ux53e3-5}

先按 \(72^\circ\)
扇区写折叠公式，再依据观察者与虚像的角差选择展开方向。分界线是
\(\theta=\phi+180^\circ\)。在这条线上，连线穿过镜面顶点，几何光学路径不成立。

\subsection{五、解析}\label{ux4e94ux89e3ux6790-5}

\textbf{单次反射。} 上镜方向为 \(36^\circ\)，关于它反射得到 \[
\theta=72^\circ-\varphi.
\] 所以第一扇区 \[
\boxed{\varphi=72^\circ-\theta,
\quad36^\circ<\theta<108^\circ}.
\] 原标尺的 \(+36^\circ\) 端对应第一扇区的 \(36^\circ\) 端，原
\(-36^\circ\) 端对应 \(108^\circ\) 端，刻度方向反转。

下镜方向为 \(-36^\circ\)，单次反射给出有向角 \[
\theta_-=-72^\circ-\varphi.
\] 将其写在 \([0,360^\circ)\)
中，\(\theta=\theta_-+360^\circ\)，所以第四扇区 \[
\boxed{\varphi=288^\circ-\theta,
\quad252^\circ<\theta<324^\circ}.
\] 这两区与观察者的角差总能沿相应单镜展开满足小于
\(180^\circ\)，因此答案不需要 \(\phi\) 分段。

\begin{figure}[H]
\centering
\includegraphics[width=0.88\linewidth,height=\textheight,keepaspectratio,alt={图6-2　正、负展开支路与对径分界}]{../figures/solution-06.pdf}
\caption*{图6-2　正、负展开支路与对径分界}
\end{figure}

\textbf{连续扇区的折叠规则。} 把展开扇区编号为整数 \(k\)，第 \(k\)
区的中心方向为 \(72^\circ k\)，角区间为 \[
72^\circ k-36^\circ<\Theta<72^\circ k+36^\circ.
\] 折回原标尺的刻度为 \[
\boxed{\varphi=(-1)^k(\Theta-72^\circ k)}.
\]
偶数次反射保持刻度方向，奇数次反射反转。沿正向展开的第二、第三区分别给出
\[
\varphi_{+}^{(2)}=\theta-144^\circ,
\qquad
\varphi_{+}^{(3)}=216^\circ-\theta.
\] 沿负向展开时使用 \(\Theta=\theta-360^\circ\)。第二扇区对应
\(k=-3\)，第三扇区对应 \(k=-2\)，因此 \[
\varphi_{-}^{(2)}=144^\circ-\theta,
\qquad
\varphi_{-}^{(3)}=\theta-216^\circ.
\]
同一个周向扇区由不同次序的镜面展开，可以给出不同的刻度方向。哪条光路可见，要由观察者位置判断。

\textbf{可见性的几何条件。} 设观察者和虚像的径向距离分别为
\(R_o,R_i>0\)。把展开方向选成 \(\Theta\)，两端角差为
\(\delta=\Theta-\phi\)。只要 \[
|\delta|<180^\circ,
\]
两端的直线段不穿原点，沿线的极角连续扫过两端之间较短的扇区，依次穿过相应展开镜面。每次交点都位于镜面的正向半直线上，折回后形成有效反射路径。

若试图沿角跨度超过 \(180^\circ\)
的展开链连线，实际直线走的是另一侧的短路径，不能依次穿过所假定的正向镜面。它对应的是把镜面无限延长之后得到的形式像，不能用来回答半无限镜的可见问题。

由于 \(36^\circ<\theta<324^\circ\) 且
\(|\phi|<36^\circ\)，正向展开成立当且仅当 \[
\theta-\phi<180^\circ,
\quad\text{即}\quad\theta<180^\circ+\phi.
\] 负向展开成立当且仅当 \[
\theta-360^\circ-\phi>-180^\circ,
\quad\text{即}\quad\theta>180^\circ+\phi.
\] 这便给出分段阈值。

\textbf{第二扇区的答案。} 对 \(108^\circ<\theta<180^\circ\)： \[
\boxed{
\varphi(\theta,\phi)=
\begin{cases}
\theta-144^\circ,&\theta<180^\circ+\phi,\\
144^\circ-\theta,&\theta>180^\circ+\phi.
\end{cases}}
\] 如果
\(\phi\ge0\)，阈值不在第二扇区内部，整个第二扇区都采用第一支。如果
\(\phi<0\)，阈值位于 \(144^\circ\) 与 \(180^\circ\)
之间，两支都可能出现。

\textbf{第三扇区的答案。} 对 \(180^\circ<\theta<252^\circ\)： \[
\boxed{
\varphi(\theta,\phi)=
\begin{cases}
216^\circ-\theta,&\theta<180^\circ+\phi,\\
\theta-216^\circ,&\theta>180^\circ+\phi.
\end{cases}}
\] 如果 \(\phi\le0\)，整个第三扇区采用第二支；如果
\(\phi>0\)，阈值进入第三扇区，前后两支分别由不同展开路径给出。

\textbf{对径线上为何不补等号。} 当
\(\theta=180^\circ+\phi\)，观察者与虚像在相反方向，任意正的径向距离都使连线通过
\(O\)。镜面交点没有唯一法线，也无法按通常的逐次平面反射继续光路。因此这个方向的像应排除，不能随意把等号划给某一支。

例如 \(\phi=-18^\circ\)，分界为
\(\theta=162^\circ\)。在第二扇区中，\(\theta=150^\circ\) 对应
\(\varphi=6^\circ\)；\(\theta=170^\circ\) 对应
\(\varphi=-26^\circ\)。中间跨过162度时看到的展开路径发生切换，而不是同一物点沿镜面平滑移动。图6-2应画出这条穿过顶点的分界直线。

\textbf{角差条件的向量证明。} 取观察者位置向量
\(\boldsymbol o\)、虚像向量 \(\boldsymbol i\)。线段为 \[
\boldsymbol z(s)=(1-s)\boldsymbol o+s\boldsymbol i,
\qquad0\le s\le1.
\] 其极角沿参数的导数符号由 \[
\boldsymbol z\times\boldsymbol z'
=\boldsymbol o\times\boldsymbol i
=R_oR_i\sin\delta
\] 决定。只要选定
\(-180^\circ<\delta<180^\circ\)，极角沿较短方向单调变化，不会额外绕原点一圈。这个证明说明分界只依赖方向角，而与观察者、标尺的径向距离无关。

\textbf{原刻度范围核验。} 第一、第四区直接有
\(-36^\circ<\varphi<36^\circ\)。第二区的两支是
\(\pm(\theta-144^\circ)\)，第三区的两支是
\(\pm(\theta-216^\circ)\)，同样都落在原刻度范围内。若所得结果达到
\(72^\circ\) 或更大，通常是折叠中心或角度零点错了。

\textbf{为何不能只按 \(360^\circ/72^\circ=5\) 计数。}
五个扇区的周向拼接会使奇数次反射的刻度方向在最后一次拼接处翻转。对连续发光标尺，某些周向区域由两种不同的镜像链重叠；半无限镜的可见支路还随
\(\phi\)
切换。简单的``共有几个像''公式无法给出每个像处的刻度函数，必须保留路径信息。

\textbf{以交点坐标验证展开光路。} 设观察者为 \(\boldsymbol o\)，虚像为
\(\boldsymbol i\)，展开镜面方向为 \(\gamma\)，单位向量为 \[
\boldsymbol e_\gamma=(\cos\gamma,\sin\gamma).
\] 直线路径写为
\(\boldsymbol x(s)=\boldsymbol o+s(\boldsymbol i-\boldsymbol o)\)。与该半直线相交时，叉积为零：
\[
\boldsymbol x(s)\times\boldsymbol e_\gamma=0.
\] 得到 \[
s_\gamma=
-\frac{\boldsymbol o\times\boldsymbol e_\gamma}
{(\boldsymbol i-\boldsymbol o)\times\boldsymbol e_\gamma}.
\] 还须同时满足 \[
0<s_\gamma<1,\qquad
\rho_\gamma=\boldsymbol x(s_\gamma)\cdot\boldsymbol e_\gamma>0.
\]
前一个条件保证交点在线段中，后一个条件保证在镜面的正向部分。只有解出交点却不检验
\(\rho_\gamma\)，相当于把半无限镜当成无限直线。

\textbf{一个两次反射的可计算例。}
取观察者半径1、\(\phi=-18^\circ\)，像半径2、\(\theta=150^\circ\)。对径分界为
\(162^\circ\)，所以用正向展开，\(k=2\)，原刻度为 \[
\varphi=150^\circ-144^\circ=6^\circ.
\] 展开路依次穿过方向 \(36^\circ\) 和 \(108^\circ\) 的镜面。两端坐标为
\[
\boldsymbol o=(\cos18^\circ,-\sin18^\circ),
\] \[
\boldsymbol i=2(\cos150^\circ,\sin150^\circ).
\] 将这两个向量代入交点式，所得两个 \(s\) 均在
\((0,1)\)，且随路径顺序递增，两交点半径也均为正。折回第二扇区后，实际路径在两面镜各反射一次，最终通向原标尺的
\(6^\circ\) 刻度。

若观察者改为 \(\phi=-30^\circ\)，同一个 \(\theta=150^\circ\)
恰位于对径分界上，线段穿过 \(O\)，应排除；若改为
\(\phi=-35^\circ\)，则分界为 \(145^\circ\)，要使用负向展开： \[
\Theta=-210^\circ,\quad k=-3,\quad
\varphi=-(-210^\circ+216^\circ)=-6^\circ.
\]
周向像位置相同，观察者换到另一侧后可见的原刻度不同。这不是公式随意改符号，而是反射链由两次改为三次。

\textbf{用反射矩阵检查刻度的方向。} 关于方向 \(\gamma\)
的平面镜反射矩阵为 \[
\mathsf R_\gamma=
\begin{pmatrix}
\cos2\gamma&\sin2\gamma\\
\sin2\gamma&-\cos2\gamma
\end{pmatrix}.
\] 两次反射复合是转角 \(2(\gamma_2-\gamma_1)\) 的旋转： \[
\mathsf R_{\gamma_2}\mathsf R_{\gamma_1}
=\mathsf Q_{2(\gamma_2-\gamma_1)}.
\] 本题镜面方向相差 \(72^\circ\)，每两次反射转动
\(144^\circ\)，并保持刻度顺序；每一次单反射行列式为
\(-1\)，使刻度顺序反转。由矩阵行列式也能检查折叠式中的
\((-1)^k\)，不必完全依赖角度记忆。

\textbf{反射次数与光程。} 对展开编号 \(k\)，有效路径的反射次数为
\(|k|\)，光程为两端直线距离： \[
L_{\rm opt}=
\sqrt{R_o^2+R_i^2-2R_oR_i\cos(\Theta-\phi)}.
\] 由于 \(\Theta\) 与 \(\theta\) 只差 \(360^\circ\)
的整数倍，候选支路的形式直线长度相同；不能只用光程更短来决定支路。可见性由交点位于正向半镜以及反射次序决定。这是半无限镜题与有限镜、无限镜题的区别。

\needspace{14\baselineskip}

\textbf{边界刻度和镜面方向的处理。} 本题给的四区间均为开区间。各端点为
\(36^\circ\)、\(108^\circ\)、\(180^\circ\)、\(252^\circ\) 与
\(324^\circ\)，不必强加进函数。像恰在镜面延长方向或原标尺端点时，可能发生光路退化或相邻展开表示重复。答题应按题给开区间写结论，只有题目额外要求连续延拓时才讨论单侧极限。

对径分界 \(\theta=180^\circ+\phi\)
位于第二或第三扇区内部时，两侧刻度一般不同。例如 \(\phi=-18^\circ\)，在
\(\theta=162^\circ\) 两侧趋于 \(+18^\circ\) 与
\(-18^\circ\)，但中间方向没有通常反射路径。因此不能用``函数必须连续''的理由把两支强行拼成一个值。

\subsection{六、关键洞察}\label{ux516dux5173ux952eux6d1eux5bdf-5}

反射代数先生成候选像，光路几何再筛选可见像。观察者角度不是多余参数；它决定第二、第三扇区在哪个方向改变反射链。

同一个像的角位置可以由 \(\theta\) 或 \(\theta-360^\circ\)
表示，但展开经过的镜面次数不同。只有与实际直线角跨度一致的展开链才有效。

\subsection{七、易错点}\label{ux4e03ux6613ux9519ux70b9-5}

把 \(\theta\)
改成眼睛处的视角；只折叠角度却不检查镜面半直线；忽略负向展开；将
\(\theta=180^\circ+\phi\)
分给任一支；认为刻度方向随扇区序号一直相同；把观察者角 \(\phi\)
与物体刻度 \(\varphi\) 混写。

\subsection{八、课堂讲解建议}\label{ux516bux8bfeux5802ux8bb2ux89e3ux5efaux8bae-5}

先讲两个单镜区，再讲中间两区。图6-2中用两色分别表示正向、负向展开，选
\(\phi=-18^\circ\) 与 \(+18^\circ\)
各画一遍对径线，学生能直接看到阈值进入哪一个扇区。

检验分段式时，既检查刻度范围，又检查射线与每条镜面相交的位置。只回代反射角公式不能发现落在镜面延长线上的无效路径。

\textbf{有限长度镜面的可见性。}
原题镜面半无限，折回后每个反射点须位于镜面正向。若进一步把镜面截成长度
\(L_m\)，还要检查其距顶点的距离小于 \(L_m\)。展开直线写为
\(\boldsymbol A+s(\boldsymbol B-\boldsymbol A)\)，镜面射线为
\(u\boldsymbol e_\gamma\)。二维叉积给 \[
s=-\frac{\boldsymbol A\times\boldsymbol e_\gamma}
{(\boldsymbol B-\boldsymbol A)\times\boldsymbol e_\gamma},
\] \[
u=\frac{\boldsymbol A\times(\boldsymbol B-\boldsymbol A)}
{\boldsymbol e_\gamma\times(\boldsymbol B-\boldsymbol A)}.
\] 必须同时满足
\(0<s<1\)、\(0<u<L_m\)，并按光路顺序检查各镜面。只知道像点角位置，并不足以保证有限镜面能看到该像。

\subsection{九、变式训练}\label{ux4e5dux53d8ux5f0fux8badux7ec3-5}

\textbf{变式一：观察者在角平分线上。} 令 \(\phi=0\)，除
\(\theta=180^\circ\) 外，四区依次为 \[
\boxed{72^\circ-\theta,\quad
\theta-144^\circ,\quad
\theta-216^\circ,\quad
288^\circ-\theta}.
\]
第二、第三扇区分别取正向二次与负向二次反射。对径方向恰在它们的共同边界，依然排除。

\textbf{变式二：追踪一个刻度。} 取原刻度 \(\varphi=10^\circ\)，观察者
\(\phi=-18^\circ\)，求第二、第三扇区内可见的角位置。

第二区正向支给 \(\theta=154^\circ\)，满足
\(154^\circ<162^\circ\)，有效；负向支给 \(134^\circ\)，不满足其应有的
\(\theta>162^\circ\)，无效。第三区全部走负向支，给
\(\theta=226^\circ\)，有效。因此 \[
\boxed{\theta=154^\circ, 226^\circ}.
\] 逐支回代条件的过程，比仅解两个线性方程更重要。

\textbf{把算出像位置与证明像可见分成两项评分。} 学生先通过反射变换得出
\(\varphi\)，再通过展开连线证明镜面交点为正半径。前者只是候选像的几何关系，后者才解决半无限镜限制。答卷只列正向四扇区函数，而完全不提观察者，第二、三区应扣除分段可见性对应的分数。

\textbf{以不等径光路排除偶然性。}
观察者半径和标尺半径无需相等。只要两端半径为正，角差小于 \(180^\circ\)
的连线，其极角沿短弧连续变化；各中间镜面均被按序穿过。角差等于
\(180^\circ\)
时，无论半径比为何都经过顶点。这就是原答案分界不含半径的原因。

可以用 \(\rho_\gamma>0\) 的交点表达式作计算核验，再交换
\(R_o,R_i\)。交换两端相当于光路逆行，几何路径仍成立。若一个实现交换后改变了可见支路，通常是只按交点顺序编号，而没有按实际展开镜面方向判断。

\textbf{推广到夹角 \(\alpha\)。} 镜面方向为 \(\pm\alpha/2\)
时，展开中心为 \(k\alpha\)： \[
\varphi=(-1)^k(\Theta-k\alpha).
\] 但 \(\theta\) 的一个周向位置可能落在多个展开链中，尤其当
\(2\pi/\alpha\) 不是偶数整数时，不可简单把该式对 \(360^\circ\)
周期延拓而丢掉路径编号。有限镜还要加入交点半径小于镜长的条件，届时可见性会依赖两端的实际距离。

\needspace{14\baselineskip}

\textbf{镜顶特殊光线。}
展开图中观察者与像点相反向时，直线穿过共同顶点。理想镜顶没有唯一法线，不能用任意一侧镜面的反射法则继续；因此等号
\(\theta=180^\circ+\phi\)
要排除。若真实镜顶有圆角，光路依赖圆角形状，而题目没有给出。

在每个有效展开路径上，应依次求与正向镜射线的交点。只验证整条直线与无限延长镜面相交，会误把背侧镜面或错误反射次序算进去。本批数值核验同时要求线段参数在零一之间、镜射线参数为正，且交点沿路单调。

二分区的结论依赖观察者刻度
\(\phi\)，而非只依赖像的分区。图6-2的临界红线用于提醒这条依赖关系，不能把某一个观察位置的分界写成全局常数角。

\newpage

\section{第 7 题}\label{ux7b2c-7-ux9898}

两根光滑平行导轨相距 \(\ell\)，左端接理想电容
\(C\)。导轨区域有垂直纸面向里的匀强磁场
\(B\)。两根导体棒横跨导轨，质量、电阻分别为 \(m_1,m_2\) 和
\(r_1,r_2\)，初速度分别为
\(v_{10},v_{20}\)，均以向右为正。电容初始不带电，并假定两棒在运动中不相碰。

（1）\(m_1=m_2=m\)，\(r_1=r_2=r\) 时，求末态共同速度 \(v_{f1}\)。

（2）求从初态到最终匀速后，第二根棒相对第一根棒的位移变化 \(S_1\)。

（3）\(m_1=m/2,m_2=m,r_1=r,r_2=r/2\) 时，求末态共同速度 \(v_{f2}\)。

（4）求此时的相对位移变化 \(S_2\)。

\includegraphics[width=0.86\linewidth,height=\textheight,keepaspectratio,alt={题7　两棒、电容与磁场的电路}]{../figures/question-07.pdf}

\subsection{一、题型定位}\label{ux4e00ux9898ux578bux5b9aux4f4d-6}

两根棒和电容之间发生电磁耦合，机械动量一般不单独守恒。运动方程与电容充电方程相加后有一个守恒组合；把各棒方程对时间积分，可直接求累计相对位移，无须先求全部指数衰减过程。

\subsection{二、知识点}\label{ux4e8cux77e5ux8bc6ux70b9-6}

动生电动势
\(B\ell v\)；欧姆定律；安培力；电容充电；守恒组合；稳定末态；电流时间积分与位移积分；能量耗散。

\subsection{三、物理图像}\label{ux4e09ux7269ux7406ux56feux50cf-6}

图7-1中，棒向右运动且磁场向里，正电荷受磁力向上，故上导轨电势较高。电容电压记为上轨对下轨
\(V\)，每根棒中向上的电流记为
\(I_i\)。图7-2沿用同一电压极性和电流方向；若某棒电动势小于 \(V\)，它的
\(I_i\) 为负，磁力也相应反向。

\subsection{四、解题入口}\label{ux56dbux89e3ux9898ux5165ux53e3-6}

记 \(b=B\ell\)，写出
\(r_iI_i=bv_i-V\)、\(m_i\dot v_i=-bI_i\)、\(C\dot V=I_1+I_2\)。求末态时消去电流；求位移时先相减欧姆式，再对时间积分。

\subsection{五、解析}\label{ux4e94ux89e3ux6790-6}

\textbf{统一符号的三组方程。}
在导轨、接线电阻忽略的模型中，两根棒两端接在同一对等势导轨上，电压均为
\(V\)。正向电流定义为棒中向上，所以 \[
\boxed{I_i=\frac{bv_i-V}{r_i}},
\qquad b=B\ell,\quad i=1,2.
\] 磁场向里，向上电流的安培力向左，因此 \[
\boxed{m_i\dot v_i=-bI_i}.
\] 两根棒把电荷送入上导轨和电容上极板： \[
\boxed{C\dot V=I_1+I_2}.
\]
三式的正负号必须同时固定。单独改变电容电压定义而不改变电流关系，会破坏后续能量式。

\textbf{守恒组合。} 将两根棒的动力学相加： \[
\frac{d}{dt}(m_1v_1+m_2v_2)=-b(I_1+I_2).
\] 利用 \(C\dot V=I_1+I_2\)： \[
\frac{d}{dt}(m_1v_1+m_2v_2+bCV)=0.
\] 由于 \(V(0)=0\)， \[
\boxed{m_1v_1+m_2v_2+bCV
=m_1v_{10}+m_2v_{20}}.
\] \(bCV\)
的单位是动量。这个守恒量来自电磁耦合方程，不能称为两根棒单独的机械动量守恒；外部磁场源与支撑导轨可以交换机械动量。

\textbf{末态共同速度。} 在最终匀速状态，\(\dot v_i=0\)，故
\(I_i=0\)。由欧姆式 \[
V_f=bv_{1f}=bv_{2f}.
\] 只要 \(B\ell\ne0\)，两棒末速度相同，记为 \(v_f\)。守恒组合给出 \[
(m_1+m_2+b^2C)v_f=m_1v_{10}+m_2v_{20}.
\] 因此通式 \[
\boxed{v_f=
\frac{m_1v_{10}+m_2v_{20}}{m_1+m_2+B^2\ell^2C}}.
\]
电阻不出现在末态速度中，因为它决定达到末态的过程和耗散速率；电容决定末态储能和守恒组合。

\begin{figure}[H]
\centering
\includegraphics[width=0.88\linewidth,height=\textheight,keepaspectratio,alt={图7-2　电流、安培力与积分守恒量}]{../figures/solution-07.pdf}
\caption*{图7-2　电流、安培力与积分守恒量}
\end{figure}

\textbf{末态确实会被趋近。} 总机械动能与电容能为 \[
\mathcal E=\frac12m_1v_1^2+\frac12m_2v_2^2+\frac12CV^2.
\] 对时间求导并代入方程： \[
\begin{aligned}
\dot{\mathcal E}
&=-bv_1I_1-bv_2I_2+V(I_1+I_2)\\
&=-(bv_1-V)I_1-(bv_2-V)I_2\\
&=\boxed{-r_1I_1^2-r_2I_2^2\le0}.
\end{aligned}
\] 只要
\(r_1,r_2>0\)，有电流的运动就持续耗散，终态满足两电流均为零。能量式既说明共同匀速是耗散后的终态，也可检验安培力和充电方程的符号。

\textbf{相对位移的通式。} 定义 \[
S=[x_2(\infty)-x_1(\infty)]-[x_2(0)-x_1(0)]
=\int_0^\infty(v_2-v_1)\,dt.
\] 它是间距的变化，不是最终间距本身。将两欧姆式相减，电容电压消去： \[
b(v_2-v_1)=r_2I_2-r_1I_1.
\] 所以 \[
S=\frac1b\left(r_2\int_0^\infty I_2dt-r_1\int_0^\infty I_1dt\right).
\] 再对每根棒的运动方程积分： \[
m_i(v_f-v_{i0})=-b\int_0^\infty I_i\,dt,
\] 即 \[
\int_0^\infty I_i\,dt=\frac{m_i}{b}(v_{i0}-v_f).
\] 得到 \[
\boxed{S=
\frac{r_2m_2(v_{20}-v_f)-r_1m_1(v_{10}-v_f)}{B^2\ell^2}}.
\]
这个通式还适用于两根棒初速度有不同符号的情形，但仍须满足不碰撞、接触保持、磁场覆盖整个过程等题设条件。

\textbf{前两问：等质量、等电阻。} 代入 \(m_1=m_2=m\)： \[
\boxed{v_{f1}=\frac{m(v_{10}+v_{20})}{2m+B^2\ell^2C}}.
\] 由于 \(r_1m_1=r_2m_2=rm\)，位移通式中的 \(v_f\) 相消： \[
\boxed{S_1=\frac{rm(v_{20}-v_{10})}{B^2\ell^2}}.
\] 若 \(v_{20}>v_{10}\)，右边的第二根棒起初更快，间距增长，\(S_1>0\)；若
\(v_{20}<v_{10}\)，间距减小，\(S_1<0\)。题目所求有向相对位移应保留符号；若求大小，再取绝对值。

\textbf{后两问：不等质量与电阻。} 当 \(m_1=m/2,m_2=m\)： \[
\boxed{v_{f2}=\frac{m(v_{10}+2v_{20})}{3m+2B^2\ell^2C}}.
\] 又有 \[
r_1m_1=r\frac m2=\frac{rm}{2},\qquad
r_2m_2=\frac r2m=\frac{rm}{2}.
\] 故 \[
\boxed{S_2=\frac{rm(v_{20}-v_{10})}{2B^2\ell^2}}.
\] 虽然各棒的质量、电阻不同，乘积 \(rm\)
相同仍使电容项在累计相对位移中消失。

\textbf{直接解相对速度复核。} 如果 \(r_1m_1=r_2m_2=K\)，从运动方程 \[
\dot v_i=-\frac{b^2}{K}v_i+\frac{b}{K}V.
\] 相减得到 \[
\frac{d}{dt}(v_2-v_1)=-\frac{b^2}{K}(v_2-v_1).
\] 因此 \[
v_2-v_1=(v_{20}-v_{10})e^{-b^2t/K}.
\] 积分为 \[
S=\frac{K}{b^2}(v_{20}-v_{10}),
\] 分别取 \(K=rm\) 与 \(rm/2\)，复现
\(S_1,S_2\)。两种推导独立核验了因子2，也说明相对速度在这两个特例中不会变号。

\textbf{不碰撞条件。} 设初始间距 \(D_0=x_2(0)-x_1(0)>0\)。当
\(v_{20}\ge v_{10}\)，间距单调增加，不会碰撞；当
\(v_{20}<v_{10}\)，间距单调减小到
\(D_0+S\)。因而在两特例中，保持正间距需要 \[
\boxed{D_0+S>0}.
\]
若恰为零，距离只在渐近末态趋于零；实际棒的有限宽度会更早碰撞。原题已假定不相碰，这一条件用于检查选用的初始数据是否符合模型。

\textbf{总耗散能量。} 令 \[
P_0=m_1v_{10}+m_2v_{20},\qquad
M_e=m_1+m_2+b^2C.
\] 末态电容电压为 \(bv_f\)，所以 \[
\mathcal E_f=\frac12M_ev_f^2=\frac{P_0^2}{2M_e}.
\] 初态电容能为零，焦耳热为 \[
\boxed{Q_J=
\frac12m_1v_{10}^2+\frac12m_2v_{20}^2
-\frac{(m_1v_{10}+m_2v_{20})^2}{2(m_1+m_2+b^2C)}}.
\]
用柯西不等式可知此量非负。若两根棒初速度相同但电容未充电，仍会先发生充电、减速和发热；``初速度相同''不是电路一开始就无电流。

\textbf{一般参数下的两条衰减模态。} 定义电动势差 \[
e_i=bv_i-V=r_iI_i.
\] 对它求导： \[
\dot e_i=-\frac{b^2}{m_ir_i}e_i
-\frac1C\left(\frac{e_1}{r_1}+\frac{e_2}{r_2}\right).
\] 令 \(a_i=b^2/(m_ir_i)\)、\(c_i=1/(Cr_i)\)，则 \[
\frac d{dt}\binom{e_1}{e_2}
=-\mathsf A\binom{e_1}{e_2},\qquad
\mathsf A=
\begin{pmatrix}
a_1+c_1&c_2\\
c_1&a_2+c_2
\end{pmatrix}.
\] 衰减率为 \[
\lambda_\pm=
\frac{\operatorname{tr}\mathsf A
\pm\sqrt{(a_1+c_1-a_2-c_2)^2+4c_1c_2}}2.
\] 两根均为实正数，因为 \[
\operatorname{tr}\mathsf A>0,\qquad
\det\mathsf A=a_1a_2+a_1c_2+a_2c_1>0.
\]
因此每个电动势差都趋于零。能量下降式说明耗散方向，正特征值进一步给出收敛速度，二者互相核对。

若初电容电压为 \(V_0\)，初值是 \(e_i(0)=bv_{i0}-V_0\)。矩阵解为 \[
\boldsymbol e(t)=e^{-\mathsf A t}\boldsymbol e(0).
\] 再积分棒方程： \[
v_i(t)=v_{i0}-\frac b{m_ir_i}\int_0^te_i(s)\,ds.
\] 这给出了完整运动，但求 \(v_f,S\)
时用前面的守恒量和积分式更简洁。完整解适用于需要某一瞬间的速度、电压或碰撞时间的追加问题。

\textbf{等质量、等电阻的显式解。} 记平均速度与速度差 \[
\bar v=(v_1+v_2)/2,\qquad w=v_2-v_1.
\] 相对模态独立： \[
\dot w=-a w,\qquad a=b^2/(mr).
\] 故 \(w=w_0e^{-at}\)。共同电动势差 \(e=b\bar v-V\) 满足 \[
\dot e=-\lambda e,\qquad
\lambda=\frac{b^2}{mr}+\frac2{Cr}.
\] 对初始 \(V_0=0\)， \[
e(t)=b\bar v_0e^{-\lambda t}.
\] 平均速度为 \[
\bar v(t)=v_f+(\bar v_0-v_f)e^{-\lambda t},
\] 其中 \[
v_f=\frac{2m\bar v_0}{2m+b^2C}.
\] 因此 \[
v_1=\bar v-\frac12w_0e^{-at},\qquad
v_2=\bar v+\frac12w_0e^{-at}.
\] 电容电压为 \[
V(t)=bv_f(1-e^{-\lambda t}).
\] 把 \(t=0\) 和 \(t\to\infty\) 代入，分别得到原初速度和共同末速度；同时
\(V(0)=0,V_f=bv_f\)。

\textbf{两种模态各自消耗多少能量。} 两棒初始动能分解为 \[
K_0=m\bar v_0^2+\frac m4w_0^2.
\] 相对模态最终消失，其能量 \(mw_0^2/4\) 全部发热。共同模态的初能为
\(m\bar v_0^2\)，最终共同动能加电容能为 \[
\left(m+\frac12b^2C\right)v_f^2.
\] 所以共同模态焦耳热为 \[
Q_{\rm common}
=\frac{mb^2C}{2m+b^2C}\bar v_0^2.
\] 总焦耳热 \[
\boxed{Q_J=\frac m4w_0^2+
\frac{mb^2C}{2m+b^2C}\bar v_0^2}.
\]
直接把通式展开也得到同一结果。电容最终储能不是总损失机械能的一半；这一比例取决于
\(b^2C/(2m)\)。

\textbf{碰撞时刻的有限时间表达。} 等 \(mr\) 的相对模态给 \[
D(t)=D_0+\frac{w_0}{a}(1-e^{-at}).
\] 若 \(w_0<0\) 且 \(D_0<-w_0/a\)，存在有限碰撞时刻： \[
t_{\rm col}=
-\frac1a\ln\left(1+\frac{aD_0}{w_0}\right).
\] 若
\(D_0=-w_0/a\)，间距只在无穷远时趋于零，没有有限时间碰撞；但它也不满足题设``保持正的末态间距''的强条件。若加入棒本身宽度，\(D\)
应改为两棒最近边缘的间距，不能继续用中心距离零作为碰撞判据。

\textbf{电容极限的次序。} \(C\to0\) 时共同模态充电很快，最终
\(v_f\to\bar v_0\)；\(C\to\infty\)
时有限时间内电压变化很慢，两棒更接近分别受电磁阻尼，最终
\(v_f\to0\)。这些极限均需保持 \(b\ne0,r_i>0\)。若直接令
\(b=0\)，两棒各自保持初速度，通常永远不同速，不能把含 \(b^{-2}\)
的位移公式代入零磁场。

\needspace{14\baselineskip}

\textbf{两根等参量棒的完整时间解。} 取 \(m_1=m_2=m,r_1=r_2=r\)，记
\(s=v_1+v_2,d=v_1-v_2\)，有 \[
\dot d=-\frac{b^2}{mr}d.
\] 不变量 \(ms+bCV=P_0=m(v_{10}+v_{20})\)。于是 \[
\dot V=\frac{bs-2V}{Cr}
=\frac{bP_0}{mCr}
-\left(\frac{b^2}{mr}+\frac2{Cr}\right)V.
\] 令 \(\lambda_s=b^2/(mr)+2/(Cr)\)，零初始电压给 \[
V(t)=b v_f(1-e^{-\lambda_s t}),
\qquad v_f=\frac{m(v_{10}+v_{20})}{2m+b^2C}.
\] 回代 \[
s(t)=2v_f+\frac{b^2C}{m}v_fe^{-\lambda_s t},
\qquad
d(t)=(v_{10}-v_{20})e^{-b^2t/(mr)}.
\] 因此每根棒速度分别为
\((s+d)/2,(s-d)/2\)，具有两个不同衰减率；只用一个指数描写各棒速度通常不正确。

积分相对位移 \[
\int_0^\infty(v_2-v_1)\,dt
=-\frac{mr}{b^2}(v_{10}-v_{20}).
\] 给定初始间距 \(L_0\)，如果前棒初速较小，累计间距减少量为
\(mr(v_{10}-v_{20})/b^2\)。它严格小于 \(L_0\)
才能避免有限时间相撞；等号时只是无限时间渐近接近，不是有限时刻碰撞。

\subsection{六、关键洞察}\label{ux516dux5173ux952eux6d1eux5bdf-6}

电容的共同电压使两根棒在末态必须有相同电动势，进而有相同速度。求终态与累计位移，不一定需要完整时间解；分别寻找可积分的方程组合，就能绕开较长的指数表达式。

累计相对位移是否与电容无关，取决于
\(r_1m_1=r_2m_2\)。它是特例中的消项条件，不是任意两棒系统的普遍性质。

\subsection{七、易错点}\label{ux4e03ux6613ux9519ux70b9-6}

直接令两棒机械动量守恒；把电容能写成
\(CV\)；两棒电流都画成固定正向却不允许负值；把 \(S\)
当作最终间距；漏掉电阻质量乘积；初电容未充电却把两棒同速当作无电流状态；令
\(B=0\) 后仍使用含 \(B^{-2}\) 的终态位移公式。

\subsection{八、课堂讲解建议}\label{ux516bux8bfeux5802ux8bb2ux89e3ux5efaux8bae-6}

在图7-1上用右手规则确定电动势极性，再画电流和安培力。先推导能量下降式检验符号，再寻找守恒组合，学生更容易理解为什么机械动量变化。

求相对位移时让学生看到``先相减，后积分''的次序。等 \(rm\)
特例可补一条指数曲线，说明位移积分收敛，但共同位置仍因最终匀速而随时间增长。

\textbf{耗散矩阵的能量检查。} 对两个杆速构成的向量
\(\boldsymbol v\)，质量矩阵为
\(M=\operatorname{diag}(m_a,m_b)\)，电磁阻尼部分可写成 \[
D_B=\frac{b^2}{R}
\begin{pmatrix}1&-1\\-1&1\end{pmatrix}.
\] 它的耗散率是 \[
\boldsymbol v^{\mathsf T}D_B\boldsymbol v
=\frac{b^2}{R}(v_a-v_b)^2=I^2R.
\] 若另有杆的线性阻力，往 \(D\) 的相应对角元加正数。无外力时 \[
\dot E=-\boldsymbol v^{\mathsf T}D\boldsymbol v\le0.
\]
电磁阻尼只衰减相对运动，不衰减共同速度；要使共同速度也趋零，必须有额外外部阻力。可用这一零模检查两杆方程的电磁力符号是否互为反向。

\subsection{九、变式训练}\label{ux4e5dux53d8ux5f0fux8badux7ec3-6}

\textbf{变式一：电容初始有电压。} 初电压为 \(V_0\)，其余条件不变。 \[
\boxed{v_f=
\frac{m_1v_{10}+m_2v_{20}+bCV_0}{m_1+m_2+b^2C}}.
\] 相对位移通式仍然成立；如果 \(r_1m_1=r_2m_2\)，\(V_0\) 不进入
\(S\)。但它会改变共同末速度和总耗散能量，不能据此说初始电容对运动没有影响。

\textbf{变式二：一根棒先向左。} 在等质量、等电阻情况下，取
\(v_{10}=-v\)、\(v_{20}=v\)。 \[
\boxed{v_f=0},\qquad\boxed{V_f=0},\qquad
\boxed{S=\frac{2rmv}{B^2\ell^2}}.
\] 两棒最终静止，间距增加，初始动能 \(mv^2\)
全部转为焦耳热。此例同时检查有向速度、最终电压与能量守恒账。

\textbf{不同 \(mr\) 时的碰撞判据。} 一般两模态解中的 \(v_2-v_1\)
是两个指数的线性组合，可能变号。因此末态间距 \(D_0+S>0\)
未必足以保证过程中不碰撞。应求 \[
D_{\min}=D_0+\min_{t\ge0}\int_0^t(v_2-v_1)\,ds.
\] 极小值候选来自 \(v_2=v_1\) 的时刻和端点。原题两个参数组恰有
\(m_1r_1=m_2r_2\)，相对速度为单指数，故前文的单调判据成立。此条件应明确说出，不能把结论提升为任意参数的通式。

\textbf{把矩阵实谱与耗散结构联系起来。} \(\mathsf A\)
虽通常不是对称矩阵，但令 \(e_i=\sqrt{r_i}\,z_i\) 后，非对角元都变为
\(1/(C\sqrt{r_1r_2})\)。它与一个实对称正定矩阵相似，因此不会出现阻尼振荡的复频率。这与忽略电感、只含电阻电容的题设一致。

若加入回路自感，电流不能瞬时满足简单欧姆式，可能产生振荡。届时原守恒组合、能量式及终态需要重新从含电感的方程推导，不能只在
\(\lambda\)
中加一个常数。用此变式可检查学生是否把电路拓扑和近似条件真正写进模型。

\needspace{14\baselineskip}

\textbf{电容初始带电的迁移。} 若初态电压为 \(V_0\)，守恒组合变为 \[
P_0=m_1v_{10}+m_2v_{20}+bCV_0.
\] 终态仍满足 \(V_f=bv_f\)，所以 \[
v_f=P_0/(m_1+m_2+b^2C).
\] 若 \(V_0\)
足够大，某棒初始电流可为负，电容先推动棒加速；总能量下降式仍为
\(-\sum r_iI_i^2\)。终态速度增加并不违背耗散，因为电容初始储能也参与了过程。

相对位移通式应从各棒积分动力学求，不宜只凭最终速度判断相撞。最终两棒同速只说明距离停止变化，不能说明此前距离一直为正。

\newpage

\section{第 8 题}\label{ux7b2c-8-ux9898}

轻质刚性圆桌的桌面为单位半径的圆，在圆周 \(N\) 等分点处装有 \(N\)
条等长桌腿。桌腿能被压缩、不能被拉伸，可视为劲度系数为 \(k\)
的单向弹簧。在桌面放置重力为 \(G\) 的质点，并取小形变近似 \(G\ll kr\)。

（1）\(N=4\) 时，求使恰好一条桌腿悬空的质点位置区域总面积 \(S_1\)。

（2）\(N\ge5\)、\(m\le N-3\) 时，求使恰好 \(m\) 条桌腿悬空的区域总面积
\(S_2\)。计算参数 \(N=21,m=15\)。

原卷前文写单位半径，形变条件中又使用 \(r\)。以下先以一般桌面半径 \(r\)
推导，原题结果令 \(r=1\)。

\includegraphics[width=0.88\linewidth,height=\textheight,keepaspectratio,alt={题8　单向弹簧腿与载荷位置}]{../figures/question-08.pdf}

\subsection{一、题型定位}\label{ux4e00ux9898ux578bux5b9aux4f4d-7}

这是刚体静力学与单向弹簧接触问题。桌面始终是平面，所以各腿的``拟压缩量''随位置线性变化；但负压缩量对应悬空，不能保留为拉力。载荷位置由非负支撑力的加权平均确定，面积问题可转化为固定接触组合对应的凸多边形面积。

\subsection{二、知识点}\label{ux4e8cux77e5ux8bc6ux70b9-7}

刚体平面形变；弹簧只承压的约束；竖直力与两个倾覆力矩平衡；连续接触块；三角函数求和；半平面交集；多边形面积；力加权重心。

\subsection{三、物理图像}\label{ux4e09ux7269ux7406ux56feux50cf-7}

图8-1用实心点表示承压桌腿，空心点表示悬空桌腿。一个平面与同一圆周上的等高点相交，正、负高度差各组成连续圆弧，因此承压腿和悬空腿各自成连续的一块，不能任意从
\(N\) 条腿中组合选择。

图8-2展示固定一块承压腿时载荷能落入的区域。把它绕圆心依次旋转
\(2\pi/N\)，得到全部等价区域；这些区域的内部对应不同的实际接触组合，不重叠。

\subsection{四、解题入口}\label{ux56dbux89e3ux9898ux5165ux53e3-7}

将承压腿数记为 \(p=N-m\)。先选取一块关于 \(x\) 轴对称的连续 \(p\)
条腿，写出各腿拟受力，再用总竖直力消去平面整体下沉量。剩余两个斜率参数受到四个边缘腿的限制，映射到载荷平面后形成一个四边形；最后乘以
\(N\)。

\subsection{五、解析}\label{ux4e94ux89e3ux6790-7}

\textbf{单向弹簧与平面条件。} 腿的周向角为 \(\vartheta_i\)，位置为 \[
\boldsymbol r_i=r(\cos\vartheta_i,\sin\vartheta_i).
\] 小倾斜下，刚性桌面对第 \(i\) 条腿的拟压缩量可写为 \[
\delta_i=a+b_x\cos\vartheta_i+b_y\sin\vartheta_i.
\] 实际支撑力为 \[
\boxed{F_i=k\max(\delta_i,0)}.
\] 负 \(\delta_i\)
表示桌面相对未压缩腿端抬高，该腿悬空，实际力为零。不能把 \(k\delta_i<0\)
当成腿向下拉桌子的力。

若载荷位于 \((X,Y)\)，静力平衡为 \[
\sum_iF_i=G,
\] \[
\sum_iF_i r\cos\vartheta_i=GX,
\qquad
\sum_iF_i r\sin\vartheta_i=GY.
\] 所以载荷位置是承压腿位置的非负力权重平均： \[
\boxed{\frac{(X,Y)}r=
\sum_i\frac{F_i}{G}(\cos\vartheta_i,\sin\vartheta_i)}.
\]
它必在承压腿的凸包内。若载荷在全部腿形成的正多边形外，仅靠非负竖直支撑力便无法保持静止；那是倾覆区，不应算成另一个正常悬空区域。

\textbf{为什么接触腿连续。} 把斜率写成
\(b\cos(\vartheta-\psi)\)，则拟压缩量为 \[
\delta(\vartheta)=a+b\cos(\vartheta-\psi).
\] 在一周上，它的正值区是围绕 \(\psi\)
的连续圆弧。离散的承压腿因此连续排列；负值腿也连续排列。在一般载荷位置下，接触组合由这个正值弧唯一确定，切换边界上才有一条或两条腿的力恰好为零。

\begin{figure}[H]
\centering
\includegraphics[width=0.9\linewidth,height=\textheight,keepaspectratio,alt={图8-2　四腿三角区域与21腿六接触区域}]{../figures/solution-08.pdf}
\caption*{图8-2　四腿三角区域与21腿六接触区域}
\end{figure}

\textbf{四腿题的直接解。} 取四腿位于
\((\pm r,0)\)、\((0,\pm r)\)。假定左腿 \((-r,0)\)
悬空，另外三腿承压。分别记右、上、下腿的力为 \(F_R,F_U,F_D\)。平衡给出
\[
F_R=\frac{GX}{r},
\] \[
F_U+F_D=G\left(1-\frac Xr\right),
\qquad F_U-F_D=\frac{GY}{r}.
\] 因此 \[
F_U=\frac G2\left(1-\frac Xr+\frac Yr\right),
\qquad
F_D=\frac G2\left(1-\frac Xr-\frac Yr\right).
\] 承压条件要求 \[
X>0,\qquad |Y|<r-X.
\] 桌面拟压缩量在四个点满足 \[
\delta_L+\delta_R=\delta_U+\delta_D.
\] 因为承压腿 \(k\delta=F\)，左腿的拟受力是 \[
k\delta_L=F_U+F_D-F_R
=G\left(1-\frac{2X}{r}\right).
\] 左腿悬空要求 \(\delta_L<0\)，故 \[
X>\frac r2.
\] 由此固定左腿悬空的区域是顶点为 \[
(r,0),\qquad(r/2,r/2),\qquad(r/2,-r/2)
\] 的三角形内部，面积 \[
A_{4,3}=\frac12\times r\times\frac r2=\frac{r^2}{4}.
\] 四条腿可分别悬空，四块三角形仅在边界接触，故 \[
\boxed{S_1=r^2}.
\] 原题单位半径时 \[
\boxed{S_1=1}.
\] 零力切换线的面积为零，是否将它纳入``悬空''的边界约定不改变面积答案。

\textbf{一般连续接触块的角度。} 以下考虑 \(1\le m\le N-3\)，因此
\(3\le p=N-m\le N-1\)。记 \[
\Delta=\frac{2\pi}{N},\qquad
\vartheta_j=\left(j-\frac{p-1}{2}\right)\Delta,\quad j=0,\ldots,p-1.
\] 这 \(p\) 条承压腿关于 \(x\) 轴对称，两端角为
\(\pm A\)，相邻两条悬空腿角为 \(\pm B\)，其中 \[
A=\frac{p-1}{2}\Delta,\qquad
B=\frac{p+1}{2}\Delta.
\] 对称性给 \(\sum\sin\vartheta_j=0\)。定义 \[
C=\sum_{j=0}^{p-1}\cos\vartheta_j
=\frac{\sin(p\Delta/2)}{\sin(\Delta/2)},
\] \[
D=\sum_{j=0}^{p-1}\cos^2\vartheta_j-\frac{C^2}{p},
\qquad
E=\sum_{j=0}^{p-1}\sin^2\vartheta_j.
\] 三角函数求和还给出 \[
\sum\cos2\vartheta_j=\frac{\sin(p\Delta)}{\sin\Delta},
\] 所以 \[
\boxed{D=\frac p2+\frac{\sin(p\Delta)}{2\sin\Delta}-\frac{C^2}{p}},
\] \[
\boxed{E=\frac p2-\frac{\sin(p\Delta)}{2\sin\Delta}}.
\] \(D,E\) 都是无量纲数；对至少三条不同位置的承压腿，二者均为正。

\textbf{用总力消去整体下沉量。} 将拟受力归一化，写成 \[
\frac{k\delta_i}{G}=a_0+u\cos\vartheta_i+v\sin\vartheta_i.
\] 对承压块求和必须为1： \[
pa_0+Cu=1.
\] 故 \[
a_0=\frac{1-Cu}{p}.
\] 任意腿的拟受力比例为 \[
\boxed{f(\vartheta)=\frac1p+
\left(\cos\vartheta-\frac Cp\right)u+v\sin\vartheta}.
\] 承压腿要求 \(f>0\)，悬空腿要求
\(f<0\)。下面求区域闭包，允许等号以便计算顶点，最终面积不变。

\textbf{四条边缘腿给出参数多边形。} 两条最外承压腿在
\(\pm A\)，两条最近悬空腿在
\(\pm B\)。满足它们的符号条件后，位于正弧内部的其他承压腿、负弧内部的其他悬空腿也满足相应条件。于是
\[
|v|\sin A\le\frac{1+(p\cos A-C)u}{p},
\] \[
|v|\sin B\le\frac{(C-p\cos B)u-1}{p}.
\] 其中 \(0<A<B\le\pi\)。当 \(B=\pi\)，第二个不等式只有对 \(u\)
的限制，按极限使用即可。

在 \(v=0\) 轴上，两端顶点为 \[
\boxed{u_a=\frac1{C-p\cos A}},
\qquad
\boxed{u_b=\frac1{C-p\cos B}}.
\] 有 \(u_a>u_b>0\)。上、下两顶点位于两类边界相交处。联立边界得到 \[
\boxed{u_d=
\frac{\sin A+\sin B}
{C(\sin A+\sin B)-p\sin(A+B)}},
\] \[
\boxed{v_d=
\frac{1+(p\cos A-C)u_d}{p\sin A}}.
\] 取 \(v_d\ge0\)，参数区域的四顶点为 \[
(u_a,0),\quad(u_d,v_d),\quad(u_b,0),\quad(u_d,-v_d).
\] 若
\(p=N-1\)，两条边界在相同的悬空腿上合并，四边形退化为三角形，上式仍可连续使用。

\textbf{参数到载荷位置的映射。} 对承压块，用 \(f_j\) 作为力权重： \[
\frac Xr=\sum_j f_j\cos\vartheta_j,
\qquad
\frac Yr=\sum_j f_j\sin\vartheta_j.
\] 由对称性 \(\sum\sin\vartheta_j\cos\vartheta_j=0\)，得 \[
\boxed{\frac Xr=\frac Cp+Du},
\qquad
\boxed{\frac Yr=Ev}.
\] 这是线性映射加平移，所以四边形变换后的两个对角线仍分别沿 \(x\) 与
\(y\)。横向对角线长为 \[
rD(u_a-u_b),
\] 竖向对角线长为 \(2rEv_d\)。固定一块 \(p\) 条承压腿的面积为 \[
A_{N,p}=\frac12[rD(u_a-u_b)](2rEv_d)
=r^2DE(u_a-u_b)v_d.
\] 因此总面积 \[
\boxed{S_2=N r^2DE(u_a-u_b)v_d,\quad p=N-m}.
\] 这个表达式由有限个三角函数组成，是一般 \(N,m\)
的解析结果。计算时先求承压数和边缘角，再求三角和，最后求参数多边形顶点；能避免把悬空腿数误作承压腿数。

\textbf{\(N=21,m=15\) 的数值。} 承压腿数为 \(p=6\)，故 \[
\Delta=\frac{2\pi}{21},\qquad A=\frac{5\pi}{21},\qquad B=\frac{7\pi}{21}.
\] 用上式计算，固定承压块的载荷区域四顶点（已除以 \(r\)）约为 \[
(0.95400983,0),\quad
(0.92047785,0.13873971),
\] \[
(0.90436759,0),\quad
(0.92047785,-0.13873971).
\] 固定区域面积为 \[
A_{21,6}\simeq0.00688734942\,r^2.
\] 21个旋转区域给出 \[
\boxed{S_2\simeq0.144634338\,r^2}.
\] 单位半径时 \(S_2\simeq0.144634338\)，占整个圆桌面积的比例为 \[
\boxed{\frac{S_2}{\pi r^2}\simeq0.0460386}.
\]
图8-2的区域应位于接近周缘的一带，不能把这类高度倾斜的载荷区画在圆心附近。

\textbf{面积为何不依赖 \(G\) 与 \(k\)。}
归一化后，支撑力之和固定为1，载荷坐标由比例权重确定；\(G/k\)
只把整个压缩量平面等比例缩放，不改变哪些腿的压缩量为正。故在小形变模型中，面积只依赖圆周几何与接触腿数。

这不意味着可任意加大载荷。小倾斜是线性压缩平面成立的条件；固定 \(N,p\)
后，若 \(G/(kr)\) 不够小，应使用实际倾斜几何与竖直弹簧接触条件重新求解。

\textbf{对所有腿的直接检查。} 为独立核验面积，可对固定承压组合在
\((u,v)\) 平面写全部 \(N\) 个线性不等式： \[
f_j\ge0\quad(j\text{在承压块}),
\qquad
f_j\le0\quad(j\text{在悬空块}).
\]
逐对求边界直线交点、删除违反任一不等式的候选点，所得凸多边形与四边缘腿公式一致。再把它的顶点代入
\((X/r,Y/r)=(C/p+Du,Ev)\)，用鞋带公式求面积，可复现
\(0.144634338r^2\)。这个复算检验了接触组合，而不只是重复代入同一个面积公式。

\textbf{\(m=0\) 的补充情况。}
如果原题允许``恰好零条悬空''，所有腿承压。由完整正多边形的对称求和， \[
F_i=\frac GN\left[1+2\frac Xr\cos\vartheta_i+2\frac Yr\sin\vartheta_i\right].
\] 要求全部 \(F_i>0\)，给出内切圆半径为 \(r/2\) 的正 \(N\)
边形。其面积为 \[
\boxed{S_{m=0}=\frac{Nr^2}{4}\tan\frac\pi N}.
\] 这个接触组合只有一个，不能像 \(m\ge1\) 时那样再乘一个额外的
\(N\)。面积边界仍不受零力腿定义影响。

\textbf{支撑力为什么是平面函数，而不是任意重心权重。}
刚性板上各腿的压缩量来自整体竖直位移与两个微小转角，故只有三个独立自由度：
\[
\delta_i=a+u_xx_i+u_yy_i.
\]
三个非共线承压腿足以确定压缩量平面。若四条或更多腿承压，它们的弹力不能任意分配，必须同时满足这个平面约束。只要求总力与总力矩平衡会留下多组力，不能据此求出哪些腿悬空。

对固定承压集合 \(\mathcal A\)，可把平衡写成矩阵： \[
k\sum_{i\in\mathcal A}
\begin{pmatrix}1\\x_i\\y_i\end{pmatrix}
(a+u_xx_i+u_yy_i)
=G\begin{pmatrix}1\\X\\Y\end{pmatrix}.
\] 记 \[
\mathsf K=k\sum_{i\in\mathcal A}
\begin{pmatrix}1\\x_i\\y_i\end{pmatrix}
\begin{pmatrix}1&x_i&y_i\end{pmatrix}.
\] 当承压点非共线时，\(\mathsf K\) 正定，可唯一解出
\((a,u_x,u_y)\)。再对所有腿检查 \(\delta_i>0\) 或
\(\delta_i\le0\)，就得到该接触集合在载荷平面的区域。前面的四边形面积是这套线性方程在圆周对称排列下的简化。

\textbf{载荷点与稳定支撑多边形的关系。} 归一化弹力 \(f_i=F_i/G\ge0\)
满足 \[
\sum f_i=1,\qquad
(X,Y)=\sum f_i(x_i,y_i).
\]
因此载荷点必须位于实际承压腿的凸包内。这个条件必要，却通常不充分：还要满足刚板---等弹簧的位移协调。前面的面积区域总在该凸包内，计算出的区域若跑到支撑点凸包外，必有符号或线性映射错误。

对于四腿，一腿悬空时剩下三腿组成三角形，静力平衡直接唯一决定三条正弹力；第四腿的拟压缩量还要非正。正弹力条件与第四腿悬空条件共同裁出图8-2的区域，不能把整个支撑三角形都算作一腿悬空区。

\textbf{力的边界与面积的边界。}
参数四边形顶点处有一条或两条腿的力恰为零。严格按``悬空''要求离开地面时，应把零力、零间隙的边界另作说明；若只统计正支撑力数，则零力腿计为不承压。两种定义只改变边界归属，不改变面积，因为有限条线段的二维面积为零。

固定承压块有 \(N\)
个旋转复制区。区内部的接触集合不同，所以互不重叠；边界处可能共享零力腿。总面积乘
\(N\) 的依据是旋转复制与内部互斥，不是载荷圆周上``看起来有 \(N\)
个相同角''。全接触集合只有一个，故其面积不能再乘一次 \(N\)。

\textbf{全接触区域的直接力分配。} 对完整正 \(N\) 边形， \[
\sum x_i=\sum y_i=\sum x_iy_i=0,\qquad
\sum x_i^2=\sum y_i^2=Nr^2/2.
\] 解平衡与协调，得到 \[
\boxed{F_i=\frac GN
\left(1+\frac{2Xx_i+2Yy_i}{r^2}\right)}.
\] 每条腿承压要求 \[
X\cos\vartheta_i+Y\sin\vartheta_i>-r/2.
\] 边界是一个内切圆半径 \(r/2\) 的正 \(N\) 边形，面积为 \[
\boxed{A_{\rm all}=\frac{Nr^2}{4}\tan\frac\pi N}.
\] 例如 \(N=4\) 面积为 \(r^2\)；\(N\to\infty\) 时趋于
\(\pi r^2/4\)。腿数增多并不使全接触面积趋于整个板面积；外围载荷仍会导致另一侧弹簧卸载。

\needspace{14\baselineskip}

\textbf{三条承压腿的接触极限。} 当 \(p=3\)
时，三支撑点的凸包仍有非零面积，协调矩阵可逆。若只剩两腿承压，载荷必须落在两腿的连线上，区域二维面积为零；只剩一腿时，载荷必须在该腿位置。原题限定
\(m\le N-3\) 正是保留有面积的区域。把 \(p=2\)
直接代入四边形面积公式，可能出现退化分母，不能据此说面积无穷大。

\textbf{小形变条件的具体量。} 总载荷为 \(G\) 时，典型压缩量为 \(G/k\)
乘无量纲系数。题设 \(G\ll kr\)
使压缩远小于板半径，支持微小转角线性化。但若板自重不忽略，应把板自重和集中载荷的合力位置合成：
\[
(X_{\rm eff},Y_{\rm eff})
=\frac{G_{\rm load}}{G_{\rm load}+G_{\rm plate}}(X,Y)
\]
（板的重心位于圆心）。总载荷改变压缩尺度，有效载荷位置则向圆心缩小，接触区域按这个映射重新求出；不能在原面积上简单加一项板自重。

\subsection{六、关键洞察}\label{ux516dux5173ux952eux6d1eux5bdf-7}

先算受力，再算面积。接触集合由桌面压缩量的符号确定，载荷位置由这些力的重心确定。把负弹簧力直接代入平衡，会给出几何上很大的区域，但它包含桌腿实际无法施加的拉力。

``恰好 \(m\) 条悬空''对应连续的一块，而不是 \(\binom Nm\)
种任意组合。利用圆周上的线性平面截面，才能正确得到总区域数 \(N\)。

\subsection{七、易错点}\label{ux4e03ux6613ux9519ux70b9-7}

把 \(m\)
当成承压腿数；允许弹簧拉桌面；未核验所有悬空腿的拟压缩量；用桌面面积代替支撑多边形的稳定范围；只求一块区域而漏乘
\(N\)；把全部接触情形也乘
\(N\)；将零力边界计入面积差；单位半径结果忘记一般半径的 \(r^2\) 因子。

\subsection{八、课堂讲解建议}\label{ux516bux8bfeux5802ux8bb2ux89e3ux5efaux8bae-7}

先完整讲四腿题，其三角区域可由三条支撑力和一条拟负力直接得到，不需要一般公式。再说明一般题只是把同样的``承压、悬空''条件写成更多线性不等式。

图8-2同时标出实际承压块和对应载荷多边形，让学生追踪顶点发生了哪条腿的零力切换。数值核验时至少取区域内一点、边界一点与区域外一点，检查总力、两个力矩和所有腿的符号。

一般结果的面积并不是猜出来的环带宽度。虽然区域旋转后接近环状，单块仍是由力平衡映射得到的四边形，圆形环带近似会丢掉边缘接触切换。

\textbf{单向弹簧平衡的变分检验。} 定义桌面平面参数
\(\boldsymbol u=(a,b_x,b_y)\)，腿位置向量
\(\boldsymbol v_i=(1,\cos\vartheta_i,\sin\vartheta_i)\)，载荷向量
\(\boldsymbol w=(1,X/r,Y/r)\)。势能为 \[
\Pi(\boldsymbol u)=\frac k2\sum_i
[\max(\boldsymbol v_i\cdot\boldsymbol u,0)]^2
-G\boldsymbol w\cdot\boldsymbol u.
\] 一阶驻值就是三个静力平衡方程。固定承压组合中， \[
\nabla^2\Pi=k\sum_{\rm active}\boldsymbol v_i\boldsymbol v_i^{\mathsf T}.
\]
三条不共线承压腿使该矩阵正定，平衡具有局部唯一性；只有一、两条承压腿时，不能稳定约束全部三个平面自由度。区域面积的内部因而必须至少有三条承压腿，退化边界上的零力腿不能算作额外的正常承压支撑。

\subsection{九、变式训练}\label{ux4e5dux53d8ux5f0fux8badux7ec3-7}

\textbf{变式一：五腿中两腿悬空。} 取 \(N=5,m=2\)，有
\(p=3\)。固定承压块区域的四顶点（除以 \(r\)）约为 \[
(1,0),\quad(0.65450850,0.47552826),\quad
(0.61803399,0),\quad(0.65450850,-0.47552826).
\] 其横、竖对角线分别为 \(0.38196601r\) 与 \(0.95105652r\)，故 \[
\boxed{S=5\times\frac12\times0.38196601\times0.95105652\,r^2
\simeq0.90817816\,r^2}.
\] 一般公式与直接三腿力平衡均得到此值。

\textbf{变式二：四腿全接触区域。} 对 \(N=4\)，全部承压要求 \[
|X|<r/2,\qquad |Y|<r/2.
\] 中央正方形面积为 \(r^2\)。四块单腿悬空三角形的总面积也为
\(r^2\)，合计 \(2r^2\)，正好等于四腿形成的支撑菱形面积。圆桌其余
\(\pi r^2-2r^2\) 的位置在支撑菱形外，不能维持正常静力平衡。

\textbf{不同腿刚度的推广应回到协调矩阵。} 若第 \(i\) 腿刚度为
\(k_i\)，有 \[
\mathsf K=\sum_{i\in\mathcal A}k_i
\begin{pmatrix}1\\x_i\\y_i\end{pmatrix}
\begin{pmatrix}1&x_i&y_i\end{pmatrix}.
\]
它仍给出固定接触集合的线性映射，但旋转对称性通常消失，不能把单块面积乘
\(N\)。可挑一条腿刚度改为 \(2k\)
的四腿例，分别列四种悬空集合，用正负压缩量检验。这个变式能区分``面积来自对称复制''与``面积与刚度绝对值无关''两个不同结论。

\textbf{总面积的数值校验。}
各种接触数对应的载荷区域可构成完整稳定载荷多边形的分区。将全接触及
\(p=3,\ldots,N-1\) 的总面积相加，应等于正 \(N\) 边形支撑点凸包面积： \[
A_{\rm hull}=\frac N2r^2\sin\frac{2\pi}{N}.
\] 两腿或一腿区只占零面积边界，不影响总和。这个校验比只检查 \(N=21,p=6\)
的单个数值更强：多乘一次旋转数、漏掉全接触区域或算错重叠区域都会使总和失配。

\textbf{载荷移出凸包时为何没有稳定静力解。}
凸组合公式要求载荷在支撑点凸包内。若载荷移出，即使允许一部分弹簧悬空，也无法靠非负竖直力平衡力矩，板将倾覆。这里讨论的是静态接触分区；倾覆后的运动、撞地或结构变形不属于这套平面小形变模型。

\needspace{14\baselineskip}

\textbf{所有接触区的覆盖检验。} 对每个承压腿数
\(p\)，把所有活跃腿与悬空腿的线性不等式一起求交，得到质心的多边形区域。只使用四条猜测边界可能遗漏另一条腿提前接触的切割。本批把一般闭式面积与全体不等式多边形的面积逐项比较，再检验所有区域面积之和等于外正多边形面积。

区域边界上有腿的压力恰为零，承压腿数的分类在这些边界有约定差异；面积不受零测度边界影响。若题目问某一点恰有几腿接地，就必须说明把零力接触计为悬空还是接触，不能用面积结论替代该约定。

支承力的凸组合使质心必须在承压接点凸包内，但这只是必要条件。相同刚板的线性位移兼容还约束各腿压缩量，因此不能直接把凸包面积当作该腿数对应的可行面积。
\end{document}
